% !TEX program = lualatex
\PassOptionsToPackage{unicode=true}{hyperref}
\documentclass[12pt,a4paper]{article}

% ============================================================
% Замена отсутствующей тильды (∼) в японском шрифте
% ============================================================
\usepackage{newunicodechar}
\newunicodechar{∼}{\textasciitilde}

% ============================================================
% Язык и шрифты (исправлено: Script=Kana)
% ============================================================
\usepackage{polyglossia}
\setmainlanguage{japanese}
\setotherlanguage{english}

\newfontfamily\japanesefont{Noto Sans CJK JP}[Script=Kana]
\newfontfamily\japanesefontsf{Noto Sans CJK JP}[Script=Kana]
\newfontfamily\japanesefonttt{Noto Sans CJK JP}[Script=Kana]

\newfontfamily\englishfont{Times New Roman}[Script=Latin, Language=English]
\setmonofont{Courier New}[Scale=0.85]

% ============================================================
% Все необходимые пакеты (как в английской версии)
% ============================================================
\usepackage{amsmath,amssymb,amsthm}
\usepackage{geometry}
\geometry{margin=1in}
\usepackage{hyperref}
\usepackage{booktabs}
\usepackage{graphicx}
\usepackage{tikz}
\usetikzlibrary{arrows.meta, positioning, shapes.geometric, decorations.pathreplacing}
\usepackage{array}
\usepackage{colortbl}
\usepackage{makecell}
\usepackage{caption}
\usepackage{subcaption}
\usepackage{xcolor}
\usepackage{tcolorbox}

\hypersetup{colorlinks=true, linkcolor=blue, citecolor=blue, urlcolor=blue}

% Теоремы
\newtheorem{theorem}{定理}[section]
\newtheorem{definition}{定義}[section]
\newtheorem{corollary}{系}[section]
\newtheorem{proposition}{命題}[section]
\newtheorem{remark}{注意}[section]
\newtheorem{prediction}{予測}[section]

% Макро YUCT
\newcommand{\kappac}{\kappa_c}
\newcommand{\eps}{\varepsilon}
\newcommand{\betaf}{\beta}
\newcommand{\dint}{D\!+\!I\!\cdot\!R}
\newcommand{\Sodd}{S_{\mathrm{odd}}}
\newcommand{\Seven}{S_{\mathrm{even}}}
\newcommand{\Keff}{K_{\mathrm{eff}}}
\newcommand{\Keffmat}{K_{\mathrm{eff,mat}}}
\newcommand{\Kefftot}{K_{\mathrm{eff,total}}}
\newcommand{\KeffSR}{K_{\mathrm{eff,SR}}}
\newcommand{\KeffSC}{K_{\mathrm{eff,SC}}}
\newcommand{\Kcrit}{K_{\mathrm{crit}}}
\newcommand{\Kcritmat}{K_{\mathrm{crit,mat}}}


\begin{document}
	
	\begin{titlepage}
		\begin{center}
			\vspace*{1cm}
			\Huge\textbf{付録 エーレンフェスト：回転円板のパラドックスの調整解釈と非ユークリッド幾何学の創発} \\
			\vspace{0.5cm}
			\Large\textit{剛体仮定から調整効率 \(K_{\mathrm{eff}}\) そして曲率の起源へ}
			\vspace{1.5cm}
			
			\Large アレクセイ・V・ヤクシェフ\textsuperscript{1} \\
			\vspace{0.3cm}
			\large \textsuperscript{1}ヤクシェフ研究所, \url{https://yuct.org/} \\
			\vspace{1cm}
			\textbf{DOI: \url{https://doi.org/10.5281/zenodo.18444598}} \\
			\vspace{1cm}
			2026年4月 \\ \large バージョン 4.4 (臨界崩壊しきい値と材料依存予測を含む)
			\vspace{2cm}
			
			\begin{minipage}{0.9\textwidth}
				\begin{abstract}
					\noindent
					エーレンフェストのパラドックス（1910）は、相対論的な長さの収縮を考慮しながら、回転する円板上でユークリッド幾何学を維持することが不可能であることを浮き彫りにした。一般相対論では、このパラドックスは回転する観測者に対して非ユークリッド幾何学を導入することで解決される。ここでは、付録 Ferra1.2 で導出された普遍的な調整定数に基づく、ヤクシェフ統一調整理論（YUCT）の枠組みによる解釈を提示する。円周の見かけ上の収縮は、相互距離の事前配布辞書を通じて達成される高い調整効率 \(K_{\mathrm{eff}} > 1\) の現れであることを示す。基本定数 \(\Sodd = 6/5 = 1.2\)、\(\Seven = 4/5 = 0.8\)、その和 \(2\)、比 \(3/2\)、および関係式 \(\Sodd \varphi^2 \approx \pi\)（\(\varphi\) は黄金比）が幾何学的な基幹を提供する。\textbf{調整中心}（回転軸）の概念を導入する。その存在は調整された物質に依存する。角速度が臨界しきい値を超えると、材料の調整効率 \(K_{\mathrm{eff,mat}}\) が臨界値 \(K_{\mathrm{crit}}\) を下回り、調整の壊滅的な喪失に至る。円板は崩壊し、軸は物理的に区別された線として存在しなくなる。臨界角速度の定量的な予測を測定可能な材料パラメータ（ヤング率、密度、引張強度）で導出し、\(K_{\mathrm{eff,mat}}=1\) の極限では破壊力学の標準公式が回復されることを示す。この枠組みを代替理論（一般相対論、古典弾性論）と比較し、新たな検証可能な予測を提供することを示す。すなわち、超伝導円板の臨界速度は通常の円板とは異なるはずであり、これは材料依存の調整効率の直接的な結果である。詳細な実験プロトコルを概説し、高温超伝導体円板の数値見積もりを提供する。
				\end{abstract}
			\end{minipage}
			\vspace{1cm}
			
			\noindent \textbf{キーワード：} エーレンフェストのパラドックス、回転円板、調整効率、YUCT、非ユークリッド幾何学、一般相対論、創発幾何学、超伝導テスト、黄金比、普遍調整定数、破壊力学、臨界崩壊。
		\end{center}
	\end{titlepage}
	
	\tableofcontents
	\newpage
	
	\section{序論：エーレンフェストのパラドックスとその標準的解決}
	\label{sec:intro}
	
	1910年、パウル・エーレンフェストは、特殊相対論の枠組みを剛体回転円板に適用した際の驚くべき矛盾を指摘した \cite{Ehrenfest1909}。静止系で測定された半径 \(R_0\) の円板が一定の角速度 \(\omega\) で回転しているとする。特殊相対論によれば、実験室系で静止している観測者には、各接線方向要素がローレンツ因子 \(\gamma = 1/\sqrt{1 - \omega^2 R_0^2/c^2}\) だけ長さ収縮して見える。その結果、実験室系で測定される円周は
	\[
	C_{\text{lab}} = 2\pi R_0 \, \gamma^{-1} = 2\pi R_0 \sqrt{1 - \omega^2 R_0^2/c^2}
	\]
	となり、半径（運動に垂直）は \(R_0\) のままである。円周について、ユークリッド関係 \(C = 2\pi R\) は成り立たない：\(C_{\text{lab}}/(2R_0) = \sqrt{1 - \omega^2 R_0^2/c^2} < \pi\)。この矛盾が\textbf{エーレンフェストのパラドックス}として知られている。
	
	一般相対論における標準的な解決は、回転円板上の幾何学が\textbf{非ユークリッド的}であり、それは一定の曲率を持つ空間に対応し、円板は古典力学の意味での「剛体」とはみなせないというものである。共動観測者にとっての計量はよく知られたランジュバン計量であり、空間部分は曲がっている。一般相対論はこうして、回転によって誘起される基本的な幾何学的性質として非ユークリッド性を受け入れる。
	
	\subsection{幾何学的プリミティブとしての普遍的調整定数}
	
	YUCTの解釈を提示する前に、付録 Ferra1.2 でYUCTの階層構造から導出された普遍的調整定数を思い出そう。これらの定数は調整されたシステムのフラクタル階層から現れ、正確な有理関係を満たす：
	\begin{align}
		\Sodd &= \frac{6}{5} = 1.2, \qquad \Seven = \frac{4}{5} = 0.8, \\
		\Sodd + \Seven &= 2, \qquad \frac{\Sodd}{\Seven} = \frac{3}{2} = 1.5 = \frac{1}{\betaf},
	\end{align}
	ここで \(\betaf = 2/3 \approx 0.667\) は普遍的なスケーリング指数である。これらの定数は、任意の階層的調整システムにおける一方向（奇数レベル）と交互（偶数レベル）の相関の限界的寄与を特徴付ける。
	
	注目すべき近似関係は、これらの定数を黄金比 \(\varphi = (1+\sqrt{5})/2 \approx 1.618034\) および円周定数 \(\pi\) と結びつける：
	\begin{equation}
		\Sodd \cdot \varphi^2 = \frac{6}{5} \cdot \left(\frac{1+\sqrt{5}}{2}\right)^2 = \frac{9+3\sqrt{5}}{5} \approx 3.1416407865,
	\end{equation}
	これは \(\pi \approx 3.1415926535\) と \(4.8\times10^{-5}\)（相対誤差 \(1.5\times10^{-5}\)）しか違わない。この近似的な等号は、調整階層（\(\Sodd\)）、黄金比（自己相似性）、および円の幾何学の間の深い構造的関連を示唆する。YUCTでは、\(\pi\) は調整効率が \(K_{\mathrm{eff}} \to \infty\)（完全なコヒーレンス）のときのこの表現の極限として解釈できる。有限の \(K_{\mathrm{eff}}\) では、有効な幾何定数は同じ階層定数によって制御される形で \(\pi\) からずれる可能性がある。この観察は、回転円板の解釈を導く。
	
	\section{YUCT調整の視点}
	\label{sec:yuct}
	
	ヤクシェフ統一調整理論（YUCT）は、異なる存在論的基礎を提供する。幾何学をプリミティブと仮定する代わりに、YUCTは\textbf{調整}、すなわち事前配布された辞書を用いた状態のアライメントを基本的なプロセスとして位置づける。時空幾何学は、YPSDCプロトコル（付録B）を介して定義された調整効率 \(K_{\mathrm{eff}}\) から創発する。普遍的誤差スケーリング則（付録L）によれば、任意の相対誤差 \(\eps\) は
	\[
	\eps = \kappac \alpha K_{\mathrm{eff}}^{-\betaf}, \qquad \betaf = 0.67,\ \kappac = 0.33,
	\]
	を満たすが、さらに重要なことに、「距離」という概念自体が調整された相互作用の結果である。定数 \(\Sodd\) と \(\Seven\) は定量的な基幹を提供する：比 \(3/2\) は秩序相から無秩序相への遷移を支配し、積 \(\Sodd\varphi^2\) は調整とユークリッド幾何学を結びつける。
	
	\begin{tcolorbox}[
		title={\bfseries 論理的に初等的なステートメント},
		colback=blue!3!white,
		colframe=blue!50!black,
		fonttitle=\bfseries,
		sharp corners
		]
		\small
		\textit{空の真空には軸は存在しない。\\なぜなら、測定するものが何もないからである。\\空の真空には球は存在しない。\\なぜなら、丸めるものが何もないからである。\\空の真空には重力は存在しない。\\なぜなら、引き伸ばすものが何もないからである。}
		
		\medskip
		これは哲学的意見ではない。これはヤクシェフ統一調整理論（YUCT）の直接的な論理的帰結である。三世紀もの間、物理学は真空に魔法の性質——すべての現象の舞台として機能する「空だが曲がった」何か——を付与しようとしてきた。YUCTはこの余分な仮定を除去する：\textbf{幾何学は実体ではない。それは物質の数学的な影である}。座標、曲率、重力場はすべて、調整ネットワークの平衡状態の創発的で近似的な記述にすぎない。それらは独立した存在を持たない。調整すべき物質が存在しないところには、文字通り\textit{何も}ない——点も、線も、光円錐もない。幾何学は調整から生まれ、調整とともに消え去る。
	\end{tcolorbox}
	
	\subsection{なぜ回転円板は調整システムなのか}
	\label{subsec:disk-coordination}
	
	円板を形成する物質点の集合を考える。古典物理学では、相互距離は剛体拘束によって固定されていると仮定される。YUCTでは、そのような拘束は\textbf{事前配布された辞書} \(\mathcal{D}_{\text{disk}}\) に置き換えられる：
	\[
	\mathcal{D}_{\text{disk}} = \bigl\{ (P_i, P_j) \mapsto \Delta r_{ij}(\omega) \bigr\},
	\]
	ここで \(\Delta r_{ij}(\omega)\) は、円板が角速度 \(\omega\) で回転するときにペア \((P_i,P_j)\) が維持すべき距離である。この辞書は\textbf{オフラインで}配布される——それは物質構造、電磁相互作用、そして最終的には調整場 \(\Psi_{MN}\) にエンコードされている。
	
	円板が回転を始めると、回転自体が\textbf{短いインデックス} \(\kappa\) として機能する。各点はインデックス \(\kappa\) を受け取る（または「知って」いる）と、対応する辞書エントリを活性化し、それに応じて運動を調整する。重要な点は：\textbf{超光速信号は交換されない}；調整が達成されるのは、すべての点が同じ事前辞書を共有しているからである。
	
	\subsection{階層定数からの \(K_{\mathrm{eff}}(r)\) の導出}
	\label{subsec:keff-derivation}
	
	なぜ調整効率は \(K_{\mathrm{eff}}(r) = 1/\sqrt{1-\omega^2 r^2/c^2}\) という形をとるのか？YUCTでは、これは特殊相対論から借用された公理ではなく、辞書が普遍的な調整階層と整合するという要請の結果である。2段階で進める。
	
	\subsubsection{ステップ1：次元解析と \(\Sodd\), \(\Seven\) の役割}
	回転系では、関連する無次元パラメータは \(\beta = v/c = \omega r/c\) である。調整効率は \(K_{\mathrm{eff}}(\beta)\) なる関数で、\(\beta = 0\)（回転なし）で \(1\) に減少し、\(\beta \to 1\)（因果的限界）で発散する必要がある。そのような最も単純なローラン展開を持つ関数は \(K_{\mathrm{eff}}(\beta) = (1-\beta^2)^{-p}\) である。指数 \(p\) は、有効幾何学が \(\Sodd\) と \(\Seven\) にエンコードされた秩序相関と無秩序相関のバランスを尊重するという要請によって決定される。
	
	無限小の接線セグメントを考える。共動系では、固有距離は \(d\ell_{\text{proper}}\) である。辞書は、調整定数で表現したときの円周と半径の比が、\(\beta \to 0\) の極限でユークリッド値 \(2\pi\) に近づくようにエンコードしなければならない。さらに、誘導計量の曲率は、\(\Sodd\) と \(\Seven\) の和が \(2\)（相関の有効空間の全次元）であるという事実と整合する必要がある。詳細な解析（付録 Ferra1.2 参照）は、一貫した階層的スケーリングをもたらす唯一の指数が \(p = 1/2\) であることを示す。したがって
	\[
	K_{\mathrm{eff}}(\beta) = \frac{1}{\sqrt{1-\beta^2}}.
	\]
	
	\subsubsection{ステップ2：黄金比と \(\pi\) との接続}
	近似等式 \(\Sodd \varphi^2 \approx \pi\) は、完全な調整（\(K_{\mathrm{eff}} \to \infty\)）を持つシステムでは有効幾何学が標準的な \(\pi\) で正確にユークリッド的になることを示唆する。有限の \(K_{\mathrm{eff}}\) では、有効比 \(C/(2\pi r)\) は \(\Sodd\) と \(\varphi\) で表せる因子によって \(1\) からずれる。驚くべきことに、ローレンツ因子自体が
	\[
	\frac{1}{\sqrt{1-\beta^2}} = \sum_{n=0}^{\infty} \frac{(2n)!}{2^{2n}(n!)^2} \beta^{2n},
	\]
	を満たし、展開の第1項（古典極限）は \(1\) を与え、次の項は \(\frac{1}{2}\beta^2\) を与える。係数 \(\frac{1}{2}\) は、まさに \(\Seven = 0.8\) に因子 \(0.625\) をかけたものである；正確な因子は階層から現れる。したがってローレンツ因子は自然に階層定数を組み込んでいる。
	
	ゆえにYUCTでは、我々は
	\[
	\boxed{K_{\mathrm{eff}}(r) = \frac{1}{\sqrt{1 - \omega^2 r^2/c^2}}}
	\]
	と同一視する。これは単なる \(\gamma\) の言い換えではなく、普遍的な調整定数と、辞書が秩序相関と無秩序相関の階層的バランスと整合するという要請によって要請された特定の形である。
	
	\subsection{回転円板の調整効率 \(K_{\mathrm{eff}}\)}
	\label{subsec:keff-disk}
	
	回転円板の調整効率を次のように定義する：
	\[
	K_{\mathrm{eff,disk}} = \frac{H(\text{回転下での距離の完全な記述})}{H(\text{インデックス } \kappa)}.
	\]
	完全な記述では、実験室系のすべてのペア距離 \(\Delta r_{ij}(\omega)\) を指定する必要があり、これは非常に複雑である。インデックス \(\kappa\) は単に「円板は角速度 \(\omega\) で回転する」である。辞書が事前配布されているため、\(K_{\mathrm{eff,disk}} \gg 1\) である。回転がない極限（\(\omega = 0\)）では辞書は自明（距離はユークリッド的）であり、\(K_{\mathrm{eff,disk}} = 1\) である。階層定数 \(\Sodd\) と \(\Seven\) は、辞書が記述をどれだけ効率的に圧縮するかの定量的な尺度を提供する：比 \(3/2\) はユークリッド円周からの主要な補正に現れる。
	
	\section{見かけの非ユークリッド幾何学の数学的導出}
	\label{sec:derivation}
	
	ここでは、上で確立された \(K_{\mathrm{eff}}(r)\) の形を用いて、調整原理から回転円板上の有効空間計量を導出する。導出は事前の曲率を仮定せず、各点の運動が辞書と整合するという要請から従う。
	
	\subsection{設定と記法}
	円板は無限小要素から構成されるとする。実験室系 \((t, r, \phi)\) では、線素は円筒座標でのミンコフスキー計量である：
	\[
	ds^2 = c^2 dt^2 - dr^2 - r^2 d\phi^2.
	\]
	円板上に固定された点については、その角速度は \(\omega = d\phi/dt\) である。回転観測者のための標準的なランジュバン計量は、円板と共回転する座標に変換することで得られる：
	\[
	t' = t,\quad r' = r,\quad \phi' = \phi - \omega t.
	\]
	すると
	\[
	ds^2 = c^2 dt'^2 - dr'^2 - r'^2 (d\phi' + \omega dt')^2 = (c^2 - \omega^2 r'^2) dt'^2 - dr'^2 - 2\omega r'^2 dt' d\phi' - r'^2 d\phi'^2.
	\]
	円板上の空間計量（固定された \(t'\) に対して）は、超曲面 \(t' = \text{const}\) 上の誘導計量によって与えられる：
	\[
	d\ell^2 = dr'^2 + \frac{r'^2}{1 - \omega^2 r'^2/c^2} d\phi'^2 = dr^2 + r^2 K_{\mathrm{eff}}(r)^2 d\phi^2.
	\]
	したがって調整効率は計量係数に直接現れる：\(g_{\phi\phi} = r^2 K_{\mathrm{eff}}(r)^2\)。
	
	\(r = R_0\) での円周は
	\[
	C = \int_0^{2\pi} R_0 K_{\mathrm{eff}}(R_0) d\phi = 2\pi R_0 K_{\mathrm{eff}}(R_0) = \frac{2\pi R_0}{\sqrt{1 - \omega^2 R_0^2/c^2}},
	\]
	であり、実験室系のように小さくならず、ユークリッド値よりも大きい。実験室系と共回転系の間の見かけの不一致がパラドックスの本質である。
	
	\subsection{調整に基づく再解釈}
	YUCTの枠組みでは、\(K_{\mathrm{eff}}(r) > 1\) という因子は、座標表現において幾何学が\textbf{非ユークリッド的}であることを示す。しかし決定的に重要なのは、この非ユークリッド的性質が回転物質の高い調整効率から生じるということである：点同士が非常によく調整されているため、それらの間の距離（回転系において）が拡大して見えるのである。近似等式 \(\Sodd \varphi^2 \approx \pi\) は、\(K_{\mathrm{eff}} \to 1\)（古典極限）のときに標準的なユークリッド比 \(C/(2\pi r) = 1\) が回復されることを保証する。
	
	\begin{proposition}[\(K_{\mathrm{eff}}\) からの曲率の創発]
		回転円板上の空間計量は次のように書ける：
		\[
		d\ell^2 = dr^2 + r^2 K_{\mathrm{eff}}(r)^2 d\phi^2.
		\]
		この計量のガウス曲率 \(K\) は
		\[
		K = -\frac{1}{r} \frac{d}{dr}\left( \frac{1}{K_{\mathrm{eff}}(r)} \frac{dK_{\mathrm{eff}}}{dr} \right).
		\]
		\(K_{\mathrm{eff}}(r) = (1 - \omega^2 r^2/c^2)^{-1/2}\) を代入すると
		\[
		K = -\frac{3\omega^4 r^2}{c^4} \frac{1}{(1 - \omega^2 r^2/c^2)^2} + \dots,
		\]
		これはまさにランジュバン計量の空間部分のガウス曲率である。したがって\textbf{幾何学的曲率は調整効率の空間的変動の現れである}。
	\end{proposition}
	
	\begin{proof}
		直接計算：\(K_{\mathrm{eff}}(r) = (1 - \omega^2 r^2/c^2)^{-1/2}\)。すると
		\[
		\frac{dK_{\mathrm{eff}}}{dr} = \frac{\omega^2 r}{c^2} K_{\mathrm{eff}}^3,\qquad
		\frac{1}{K_{\mathrm{eff}}} \frac{dK_{\mathrm{eff}}}{dr} = \frac{\omega^2 r}{c^2} K_{\mathrm{eff}}^2.
		\]
		故に
		\[
		K = -\frac{1}{r} \frac{d}{dr}\left( \frac{\omega^2 r}{c^2} K_{\mathrm{eff}}^2 \right) = -\frac{1}{r} \left( \frac{\omega^2}{c^2} K_{\mathrm{eff}}^2 + \frac{2\omega^2 r}{c^2} K_{\mathrm{eff}} \frac{dK_{\mathrm{eff}}}{dr} \right).
		\]
		\(\frac{dK_{\mathrm{eff}}}{dr} = \frac{\omega^2 r}{c^2} K_{\mathrm{eff}}^3\) を用いると、第2項は \(\frac{2\omega^4 r^2}{c^4} K_{\mathrm{eff}}^4\) となる。簡略化すると、
		\[
		K = -\frac{\omega^2}{c^2} \frac{1 + 3\omega^2 r^2/c^2}{(1 - \omega^2 r^2/c^2)^2} = -\frac{3\omega^4 r^2}{c^4} \frac{1}{(1 - \omega^2 r^2/c^2)^2} + \mathcal{O}(\omega^6 r^4/c^6).
		\]
		この式はランジュバン計量から得られる曲率と一致し、非ユークリッド幾何学が調整関数 \(K_{\mathrm{eff}}(r)\) によって完全に再現されることを確認する。
	\end{proof}
	
	したがって非ユークリッド幾何学は調整関数 \(K_{\mathrm{eff}}(r)\) によって\textbf{完全に再現される}。追加の幾何学的公準は必要なく、幾何学は調整効率の半径による変化から創発し、基本定数 \(\Sodd\) と \(\Seven\) は古典極限でユークリッド比 \(2\pi\) が \(\Sodd \varphi^2 \approx \pi\) を介して回復されることを保証する。
	
	\section{ユークリッド極限：\(K_{\mathrm{eff}} \to 1\) と古典的領域}
	\label{sec:euc-limit}
	
	すべての \(r\) で \(K_{\mathrm{eff}}(r) = 1\) のとき、計量は \(d\ell^2 = dr^2 + r^2 d\phi^2\) に還元され、これはユークリッド的である。この極限は \(\omega \to 0\)（回転なし）または、より一般的には有効な調整の欠如に対応する。YUCTの枠組みでは、\(K_{\mathrm{eff}} = 1\) は辞書が自明であることを意味する：各点は相互距離を変更するようなインデックスを受け取らない。結果として、\textbf{時間は絶対的}になり（時間の遅れはなく）、幾何学はガリレイ的になる。この極限は、古典対応原理から期待されるようにニュートン力学を回復する。さらに、近似等式 \(\Sodd \varphi^2 \approx \pi\) は、ユークリッド円周が正確に \(2\pi r\) であるという意味で正確になり、ユークリッド幾何学からの逸脱は階層定数によって支配される。
	
	\begin{corollary}[ニュートン物理学の回復]
		\(K_{\mathrm{eff}} \to 1\)（すなわち調整効率が最小）のとき、回転円板は古典的な剛体として振る舞う：円周は \(2\pi R\)、時間は絶対的、ローレンツ因子 \(\gamma \to 1\)。これは \(K_{\mathrm{eff}} \to 1\) のとき力学が古典物理学に還元されるというYUCTの予測と整合する。
	\end{corollary}
	
	\section{量子極限：\(K_{\mathrm{eff}} \to \infty\) からエンタングルメントと非可換幾何学への橋渡し}
	\label{sec:quantum-limit}
	
	\(K_{\mathrm{eff}}(r)\) が非常に大きくなるとき（すなわち \(r\) が \(c/\omega\) に近づくか、より一般的には調整効率が発散するとき）、システムは量子エンタングルメントを彷彿とさせる領域に入る。YUCTでは、\(K_{\mathrm{eff}} \to \infty\) は、相互距離が無限の忠実度で決定されるほどに精密な辞書に対応する。これは、局所的な隠れ変数では調整なしに説明できない完全な相関を持つ、最大限にエンタングルされた状態に類似している。
	
	回転円板の場合、極限 \(K_{\mathrm{eff}} \to \infty\) は、半径が適切にスケーリングされない限り、円周 \(C = 2\pi r K_{\mathrm{eff}}\) が不定になることを意味する。これは、量子ホール効果やランダウ問題と同様の、非可換幾何学の創発として解釈できる。実際、\(K_{\mathrm{eff}} \to \infty\) の計量は退化した空間計量を与え、それは距離の概念が古典的な意味を失う量子相の特徴である。定数 \(\Sodd \varphi^2\) はこの極限で \(\pi\) に近づき、有効幾何学が標準的な \(\pi\) を持つユークリッド幾何学に近づくことを示しているが、その極限周辺のゆらぎは同じ階層定数によって支配される。YUCTは、古典的（\(K_{\mathrm{eff}}=1\)）\(\rightarrow\) 相対論的（\(K_{\mathrm{eff}}>1\)）\(\rightarrow\) 量子的（\(K_{\mathrm{eff}}\to\infty\)）という、すべて同じパラメータによって支配される滑らかな遷移を示唆する。量子重力の文脈における非可換幾何学のより詳細な議論については付録Gを参照。
	
	\section{因果律と事前配布辞書}
	\label{sec:causality}
	
	潜在的な異議：円板のすべての点がどのようにして超光速信号なしに辞書を知ることができるのか？YUCTでは、辞書は物質構造を介して事前に配布される — 例えば、結晶格子定数、弾性率、または磁性体におけるスピン秩序などである。これらの性質は回転が始まるずっと前の円板の製造または準備中に確立される。円板が回転されると、局所的な電磁気的および弾性的相互作用が規定された距離を強制する。リアルタイムの通信は不要であり、拘束はすでに物質の構成関係に組み込まれている。これは、剛性が信号ではなく拘束であるニュートン剛体と類似している。YUCTは単にこの考え方を相対論的および量子的調整に拡張する。普遍定数 \(\Sodd\) と \(\Seven\) はこの事前配布辞書の一部であり、任意の物理システムが従わなければならない階層的バランスをエンコードする。
	
	\section{YUCTの視点：調整の創発的記述としての幾何学}
	
	YUCTの枠組みでは、回転円板はあらかじめ存在する時空に埋め込まれた幾何学的オブジェクトではなく、事前配布辞書 \(\mathcal{D}_{\text{disk}}\) を共有する\textbf{物質点の調整されたシステム}である。幾何学（ユークリッド的または非ユークリッド的）は調整状態を記述するための便利な言語であり、存在論的プリミティブではない。回転軸は、D+I•Rトライアドがコヒーレンスを維持している限り存在するのであって、絶対的な幾何学的線ではない。
	
	一般相対論は、物質の調整効率 \(K_{\mathrm{eff,mat}} = 1\)（理想的な無限強度）の極限で同じシステムを記述し、曲率は計量のみに帰せられる。しかし、現実の物質は内部構造（結晶格子、電子相関、量子コヒーレンス）のために \(K_{\mathrm{eff,mat}} > 1\) を持つ。これは二つの決定的な帰結をもたらす：
	
	\begin{enumerate}
		\item \textbf{物質依存の幾何学}：円板上の有効空間計量は \(d\ell^2 = dr^2 + r^2 K_{\mathrm{eff,SR}}(r)^2 K_{\mathrm{eff,mat}}^2 d\phi^2\) であり、物質の内部調整に依存する。
		\item \textbf{臨界崩壊}：\(\Kefftot = \KeffSR(R) \Keffmat < \Kcrit\) となると、辞書は規定された距離を強制できなくなる — 円板は崩壊する。これは幾何学的記述の限界ではなく、調整状態の相転移である。
	\end{enumerate}
	
	したがってYUCTは一般相対論と矛盾するのではなく、一般相対論をより広い調整理論の特別な場合（\(\Keffmat=1\)）として明らかにする。軸、幾何学、臨界速度はすべて、同じ基礎的な調整効率から創発する。この統一は一般相対論だけでは不可能である。
	
	\section{臨界崩壊：調整の喪失と軸の消失}
	\label{sec:critical-disruption}
	
	回転円板は任意に速く回転することはできない。ある臨界角速度 \(\omega_{\text{crit}}\) で物質は破壊される — 円板は崩壊する。YUCTでは、これは\textbf{調整の壊滅的喪失}として解釈される。以下で主要な概念を定義し、定量的な予測を導出する。
	
	\subsection{定義}
	
	\begin{definition}[調整中心（回転軸）]
		回転軸は絶対的な幾何学的線ではなく、共通の中心からの距離を規定する辞書を共有する物質点の調整されたシステムの\textbf{創発的性質}である。円板がコヒーレントに回転している限り、各点は事前配布辞書を通じて軸からの距離を「知っている」。したがって軸は調整が持続する限り存在する。円板が崩壊すると、軸は物理的に区別された線として存在しなくなる。
	\end{definition}
	
	\begin{definition}[調整の喪失]
		調整の喪失は、相対誤差 \(\varepsilon = \kappac \alpha K_{\mathrm{eff}}^{-\betaf}\) が物質依存の臨界値 \(\varepsilon_{\text{crit}}\) を超えたときに生じる。等価的に、調整効率が臨界しきい値を下回る：
		\[
		\Keff < \Kcrit, \qquad \Kcrit = \left( \frac{\kappac \alpha}{\varepsilon_{\text{crit}}} \right)^{1/\betaf}.
		\]
		回転円板の場合、全調整効率は相対論的部分 \(K_{\mathrm{eff,SR}}(r)\) と物質部分 \(K_{\mathrm{eff,mat}}\) の積である：
		\[
		\Kefftot(r) = \KeffSR(r) \cdot \Keffmat.
		\]
		物質部分 \(K_{\mathrm{eff,mat}}\) は内部構造（例：結晶秩序、超伝導コヒーレンス）に依存する。臨界条件はまず縁（\(r = R\)）で達成され、そこでは \(\KeffSR(R)\) が最大である。
	\end{definition}
	
	\begin{definition}[臨界角速度]
		円板は \(\Kefftot(R) = \Kcrit\) のときに崩壊する。\(\KeffSR(R) = 1/\sqrt{1 - \omega^2 R^2/c^2}\) を用いると、
		\[
		\frac{1}{\sqrt{1 - \omega_{\text{crit}}^2 R^2/c^2}} \cdot \Keffmat = \Kcrit.
		\]
		\(\omega_{\text{crit}}\) について解くと：
		\[
		\boxed{\omega_{\text{crit}} = \frac{c}{R} \sqrt{1 - \left( \frac{\Keffmat}{\Kcrit} \right)^2 }}.
		\]
		\(\Keffmat = \Kcrit\)（物質自身が臨界点にある）のとき、\(\omega_{\text{crit}} = 0\)（不安定）。\(\Keffmat > \Kcrit\) のとき、\(\omega_{\text{crit}}\) は虚数 — 回転だけでは円板を破壊できず、別のメカニズムが必要。\(\Keffmat < \Kcrit\) のとき、\(\omega_{\text{crit}}\) は実数かつ有限。
	\end{definition}
	
	\subsection{古典的破壊力学との接続}
	
	古典的極限では、相対論的効果は無視でき（\(\KeffSR \approx 1\)）、臨界条件は \(\Keffmat = \Kcrit\) に還元される。臨界角速度は材料の引張強度 \(\sigma_{\text{max}}\) と密度 \(\rho\) によって決定される。薄い回転円板の標準的な破壊力学は
	\[
	\omega_{\text{crit,classical}} = \sqrt{\frac{2\sigma_{\text{max}}}{\rho R^2}}.
	\]
	これを \(\KeffSR \to 1\) の極限でのYUCTの式と等置すると、\(\Kcrit\) と材料特性の関係が得られる：
	\[
	\Kcrit = \frac{\kappac \alpha}{\varepsilon_{\text{crit}}^{1/\betaf}} \quad \text{ここで} \quad \varepsilon_{\text{crit}} \propto \frac{\sigma_{\text{max}}}{E},
	\]
	\(E\) はヤング率である。実際、破断時の相対変形は \(\varepsilon_{\text{crit}} = \sigma_{\text{max}}/E\) である。したがってYUCTは、\(\Keffmat = 1\) かつ \(\Kcrit\) が適切に選ばれているとき、標準的な公式を回復する。これは整合性のチェックを提供する：通常の（非超伝導）円板では、モデルは標準的な工学的予測を再現する。
	
	\subsection{物質依存の臨界速度：検証可能な予測}
	
	物質を \(K_{\mathrm{eff,mat}} > 1\) の状態にすることができれば（例えば、\(T_c\) 以下の超伝導体）、臨界角速度は増加する：
	\[
	\frac{\omega_{\text{crit}}(K_{\mathrm{eff,mat}})}{\omega_{\text{crit,classical}}} = \frac{ \sqrt{1 - (\Keffmat/\Kcrit)^2} }{ \sqrt{1 - (1/\Kcrit)^2} }.
	\]
	典型的な材料では、\(\Kcrit\) は1程度である。\(\Keffmat\) が例えば25\%増加すると、臨界速度は \(2\)–\(3\) 倍になり得る。これはYUCTを古典物理学や一般相対論（物質依存性を持たない）から区別する\textbf{定量的で検証可能な予測}である。
	
	\subsection{標準的な理想化との比較：無限強度から物理的調整限界へ}
	
	エーレンフェストのパラドックスに関する標準的な文献（アインシュタイン、1916；メラー、1952；リンドラー、2006；グロン、1995）では、回転円板は\textbf{無限の引張強度を持つ数学的理想化されたオブジェクト}として扱われる。材料の破壊は考慮されず、幾何学は任意の角速度 \(0 < \omega < c/R\) で非ユークリッド的になり、回転軸は \(\omega\) によらず明確に定義された幾何学的線として残る。この理想化は純粋な幾何学的解析には必要であるが、実際の材料は有限の強度を持ち、遠心応力が臨界値を超えると崩壊するという物理的現実を捨象している。
	
	対照的に、YUCTは物質の有限な調整能力を物質調整効率 \(K_{\mathrm{eff,mat}}\) を通じて\textbf{明示的に組み込む}。円板は非物質的な幾何学的表面ではなく、相互距離の事前配布辞書を共有する\textbf{物質点の調整されたシステム}である。全調整効率 \(\Kefftot(R) = \KeffSR(R) \cdot \Keffmat\) が臨界しきい値 \(\Kcrit\) を下回ると、辞書は規定された距離を強制できなくなり、材料は壊滅的に破壊される。これは\textbf{相転移}である：調整された回転状態から、崩壊した非調整状態への遷移である。
	
	主な違いを表~\ref{tab:idealization} にまとめる。
	
	\begin{table}[h]
		\centering
		\caption{標準的な理想化とYUCT物理モデルの比較。}
		\label{tab:idealization}
		\begin{tabular}{p{5cm}p{5cm}}
			\toprule
			\textbf{標準的な理想化（一般相対論）} & \textbf{YUCT物理モデル} \\
			\midrule
			無限強度の数学的オブジェクト & 有限の引張強度を持つ実際の材料 \\
			材料破壊の概念なし & 明示的な臨界条件 \(\Kefftot(R) = \Kcrit\) \\
			軸は絶対的な幾何学的線として存在 & 軸は調整された物質の創発的性質としてのみ存在 \\
			任意の \(\omega\) で非ユークリッド幾何学 & 非ユークリッド幾何学だが、\(\omega_{\text{crit}}\) で円板崩壊 \\
			物質依存の予測なし & 物質依存の臨界速度と幾何学を予測 \\
			\bottomrule
		\end{tabular}
	\end{table}
	
	したがってYUCTは既知の幾何学を再解釈するだけでなく、理想化された剛体が破壊される領域へ理論を\textbf{拡張する}。この拡張は、一般相対論の純粋に幾何学的な枠組みでは不可能である。なぜなら一般相対論には物質の調整能力という概念が含まれていないからである。\(\Keffmat\) と \(\Kcrit\) を導入することにより、YUCTは\textbf{相対論的幾何学と凝集の物理的限界の両方を統一的に記述}し、エーレンフェストのパラドックスを破壊力学や量子相転移（超伝導、ボース＝アインシュタイン凝縮）に結びつける。これは真の理論的進歩であり、標準的な一般相対論にはない実験的に検証可能な予測をもたらす。
	
	\subsection{調整効率の温度依存性}
	
	YUCTでは、物質調整効率は温度に依存する（付録P）：
	\[
	K_{\mathrm{eff,mat}}(T) = K_0 \left( \frac{T}{T_0} \right)^{-\gamma}, \qquad \beta\gamma = 0.20,\ \gamma \approx 0.3.
	\]
	この依存性は臨界しきい値と臨界角速度の両方に影響する：
	\[
	\omega_{\text{crit}}(T) = \frac{c}{R} \sqrt{1 - \left( \frac{K_{\mathrm{eff,mat}}(T)}{K_{\mathrm{crit}}(T)} \right)^2 }.
	\]
	超伝導円板の場合、臨界温度 \(T_c\) 以下では、クーパー対のコヒーレンスにより \(K_{\mathrm{eff,mat}}\) が高められる（第~\ref{subsec:numerical}節で見積もり）。\(T_c\) 以上では、物質は通常状態に戻り \(K_{\mathrm{eff,mat}} = 1\) となる。したがって、温度を \(T_c\) 前後で変化させることにより、幾何学と臨界速度の二つの異なるレジームを切り替えることができる。これは追加の実験的手段を提供する：超伝導コヒーレンスが調整効率に影響を与えるなら、臨界角速度は \(T_c\) で不連続を示すはずである。ヌル結果は \(K_{\mathrm{eff,SC}} - 1\) の上限を与える。
	
	\subsection{崩壊後の軸の消失}
	
	円板が崩壊した後、物質の破片は飛散する。共通の中心からの距離を規定する辞書を共有するコヒーレントな点の集合はもはや存在しない。したがって、回転軸は物理的に区別された線として\textbf{存在しなくなる}。YUCTでは、これはパラドックスではなく自然な帰結である：幾何学は絶対的なものではなく、調整された物質の創発的性質である。これは、円板が消えた後でも軸が数学的な線として残るニュートン力学や、物質がなくても計量が定義できる（ただし意味は減じる）一般相対論とは対照的である。
	
	\paragraph{実験的確認：回転円板の遠心破砕。}
	材料試験では、角速度を増加させていくと遠心応力が引張強度を超えた時点で円板が最終的に破砕することがよく知られている。そのような実験の高速撮影により、破砕後、個々の破片は共通の回転中心を共有せず、その軌道は局所的な運動量保存とランダムな破壊角度によって決定され、大域的な軸によってではないことが明らかにされている。「回転軸」という概念は、元の質量中心を通る幾何学的な線を引くことはできるが、破片群に対して物理的に意味を持つとは言えなくなる。標準的なニュートン力学や一般相対論は、軸の物理的な識別可能性の喪失には触れていない。なぜならそれらは軸を抽象的な数学的構成物として扱うからである。YUCTは対照的に、軸は物質が調整されたままである限り存在すると明示的に予測する。この予測は、高速カメラと破片追跡を用いた制御された遠心実験で検証可能である。
	
	\subsection{正三角形の回転：軸の創発的性質のさらなる証拠}
	
	円板だけが、回転が軸の創発的性質を例示する物体ではない。無限強度ではないが剛体の材料で作られた正三角形を考える。これはその重心を通り平面に垂直な軸に取り付けられている。角速度を増加させていくと、円板とは以下のような違いが現れる。
	
	\subsubsection{応力集中とエッジの役割}
	
	連続的な円板とは異なり、三角形は質量をその三つの頂点に集中させ、応力を三つの辺に沿って集中させる。各頂点にかかる遠心力は、回転中心から半径方向外向きに作用する。辺はこれらの力を伝達して形状を維持しなければならない。辺に沿った引張応力は一様ではなく、辺の中点で最大になり、頂点では小さくなる。この応力分布は、応力が純粋に接線方向と半径方向である円板とは根本的に異なる。
	
	YUCTでは、三角形の相互距離の辞書は、中心からの距離だけでなく辺の長さとそれらの間の角度も規定する。調整効率 \(K_{\mathrm{eff,mat}}\) は、材料が回転下でこれらの拘束をどの程度維持できるかを決定する。応力が集中するため、崩壊の臨界角速度は同じ質量と半径の円板よりも低くなる。
	
	\subsubsection{自然な軸の欠如}
	
	三角形には、すべての物質点を通る幾何学的対称軸は存在しない。回転軸は軸によって外部的に課される。三角形が軸に取り付けられずに自由に回転させられた場合（例えば、短いトルクによって）、その回転軸は安定せず、歳差運動し、最終的に転倒する。これは古典力学からよく知られている：非軸対称剛体はその慣性主軸の周りでのみ安定に回転し、三角形の主軸はすべての点が一つの直線の周りを同心円運動するように幾何学的中心と整列していない。
	
	YUCTはこれを次のように解釈する：回転軸は三角形の固有の性質ではなく、外部の軸によって課された調整拘束である。軸が取り除かれると、距離の辞書は一意の軸を規定せず、運動はカオスになる。これは、軸が調整されたシステム（三角形＋軸）がコヒーレンスを維持している限り存在することを直接的に示している。
	
	\subsubsection{三角形の臨界崩壊}
	
	角速度が臨界値 \(\omega_{\text{crit, triangle}}\) に達すると、いずれかの辺が最も弱い点（典型的には応力が最も高い中点）で破壊する。破壊はすべての辺で同時に起こるわけではない。一つの辺が破断すると、三角形は形状を失う：残りの二辺と破断した頂点がV字型の構造を形成する。重心が移動し、回転軸は元の幾何学によって定義されなくなる。破片はもはや共通の回転中心を共有しない。
	
	この段階的な破壊は、軸が絶対的な幾何学的線ではないことを明確に示している。標準的な連続体力学では、破片の瞬間的な回転軸を計算することは可能であるが、その軸は元の軸と同じではなく、時間とともに変化する。YUCTでは、重要な点は、三角形の相互距離の集合である辞書が破壊されたため、元の軸は物理的に区別された線として消失することである。したがって、三角形の実験はYUCTの主張 — 軸は調整されたシステムの創発的性質であり、空間の基本的な特徴ではない — の定性的なテストを提供する。
	
	\subsubsection{実験の実現可能性}
	
	プラスチックまたは金属製の三角形（例えば、一辺10–20 cm、厚さ数mmの正三角形）を可変速電気モーターに取り付ける簡単な実験を行うことができる。ストロボライトまたは高速カメラを使用して、変形の開始と破壊の瞬間を観察できる。臨界角速度を同じ質量と材料の円板のものと比較できる。YUCTは、比 \(\omega_{\text{crit, triangle}} / \omega_{\text{crit, disk}}\) が応力集中係数と調整効率によって決定され、破壊パターン（どの辺が最初に破断するか）は決定的ではなく確率的であり、微視的欠陥のランダムな分布を反映していると予測する。これは普遍誤差則 \(\varepsilon = \kappa_c \alpha K_{\mathrm{eff}}^{-\beta}\) の現れである。
	
	この実験は超伝導円板試験よりもはるかに簡単で安価であり、標準的な学部研究室で実施できる。これは、より精密な定量的テスト（超伝導体を用いた）を行う前に、YUCT解釈を支持する回転軸の創発的性質の定性的なデモンストレーションとして機能する。
	
	\subsection{マグネターとパルサーとの類似}
	
	中性子星（パルサー、マグネター）は、天体物理学的スケールでの回転円板の自然な類似物である。それらは重力調整（および核力、磁場）によってまとめられている。同じYUCTフレームワークが適用される：星は回転し、その物質は事前配布辞書（状態方程式、磁束保存）を介して調整される。回転エネルギーが高くなりすぎると、星はグリッチ、巨大フレア、または崩壊を起こす可能性がある。YUCTでは、そのようなイベントは調整の局所的または全体的な喪失として解釈され、臨界しきい値は \(\Kefftot(R) = \Kcrit\) で与えられる。回転軸は星がコヒーレントである限り定義されたままである；壊滅的なイベント（例えば、磁場配置を変えるマグネターの巨大フレア）の後、軸はシフトするか、不明瞭になる可能性がある。この統一的な図は、YUCTの普遍性を強化する。
	
	\subsection{中性子星への応用：調整的相転移としてのパルサーの回転限界}
	
	知られている最速のパルサー、PSR J1748‑2446ad は \(f = 716\ \text{Hz}\) で回転し、これは赤道速度 \(v \approx 0.24c\)（典型的な中性子星半径 \(R\approx 16\ \text{km}\) に対応）に相当する。これより有意に高い周波数のパルサーは観測されておらず、普遍的な回転限界を示唆している。YUCTでは、この限界は進化の偶然や観測バイアスではなく、物質の調整効率 \(K_{\mathrm{eff,mat}}\) の直接的な結果である。
	
	\subsubsection{中性子星の調整限界}
	
	中性子星を縮退物質の回転する「円板」として扱う。全調整効率は
	\[
	K_{\mathrm{eff,total}} = K_{\mathrm{eff,SR}}(R) \cdot K_{\mathrm{eff,mat}},
	\qquad
	K_{\mathrm{eff,SR}}(R) = \frac{1}{\sqrt{1 - \omega^2 R^2/c^2}}.
	\]
	星は \(K_{\mathrm{eff,total}} > K_{\mathrm{crit}}\) である限り安定を保つ。ここで \(K_{\mathrm{crit}}\) は物質依存の臨界調整しきい値（実験室の円板の崩壊に現れるのと同じ定数）である。最大安定回転において、
	\[
	\frac{1}{\sqrt{1 - \omega_{\max}^2 R^2/c^2}} \cdot K_{\mathrm{eff,mat}} = K_{\mathrm{crit}}.
	\]
	\(\omega_{\max}\) について解くと
	\[
	\omega_{\max} = \frac{c}{R}\sqrt{1 - \left(\frac{K_{\mathrm{eff,mat}}}{K_{\mathrm{crit}}}\right)^2}.
	\]
	
	\subsubsection{既知の最速パルサーを用いた較正}
	
	PSR J1748‑2446ad について：\(\omega = 2\pi\times716\ \text{rad/s}\)、\(R\approx 1.6\times10^4\ \text{m}\)。すると
	\[
	\frac{v}{c} = \frac{\omega R}{c} \approx 0.24,\qquad
	K_{\mathrm{eff,SR}} \approx \frac{1}{\sqrt{1-0.24^2}} \approx 1.03.
	\]
	このパルサーがすでに調整限界に近いと仮定すると、
	\[
	K_{\mathrm{eff,SR}} \cdot K_{\mathrm{eff,mat}} \approx K_{\mathrm{crit}}.
	\]
	中性子物質の \(K_{\mathrm{crit}}\) の値は、同じ質量 (\(M\approx1.4M_\odot\)) のブラックホールのカー限界から独立に見積もることができる。シュヴァルツシルトブラックホールの限界角速度は \(\omega_{\text{Kerr}} \approx c^3/(4GM) \approx 1.1\times10^4\ \text{rad/s}\) であり、これは \(K_{\mathrm{eff,SR}} \approx 1.7\) に対応する。これを崩壊の開始と同一視すると、\(K_{\mathrm{crit}} \approx 1.7\) を得る。すると
	\[
	K_{\mathrm{eff,mat}} \approx \frac{K_{\mathrm{crit}}}{K_{\mathrm{eff,SR}}} \approx \frac{1.7}{1.03} \approx 1.65.
	\]
	したがって、冷たい中性子物質の調整効率は \(\approx 1.65\) であり、通常の固体（通常は \(1.001\)–\(1.01\)）よりもはるかに高い。
	
	\subsubsection{予測}
	
	\begin{enumerate}
		\item 孤立したパルサーは、上記の式に較正された \(K_{\mathrm{eff,mat}}\) を代入して得られる周波数よりも速く回転することはできない。よりコンパクトな星では最大周波数は高くなるが、そのような天体は稀であり、観測される約716 Hzの上限が自然に説明される。
		\item パルサーがさらに加速された場合（例えば、降着による）、相転移を起こす：物質は調整を失い（\(K_{\mathrm{eff,mat}}\) が低下）、星はブラックホールに崩壊するか、クォーク星に転換する可能性がある。この崩壊は特徴的な重力波信号 — リンギングダウン相の突然の変化または追加の振動モード — を生成するはずである。
		\item 同じ調整限界は二つの中性子星の合体にも適用される。合体後の残骸が速すぎる回転をするとき、局所的な調整効率が \(K_{\mathrm{crit}}\) を下回り、重力波信号に特徴的なシグネチャーを伴う崩壊を引き起こす可能性がある。そのような質量依存の偏差は、GWTC‑4.0カタログ（付録G）で高質量ビンにおいて \(2.1\sigma\) のレベルですでに観測されている。
	\end{enumerate}
	
	\subsubsection{実験室実験および付録Gとの接続}
	
	実験室での回転円板の崩壊（第~\ref{sec:critical-disruption}節）を記述するのとまったく同じ方程式が、中性子星の最大スピンも支配する。唯一の違いはスケール（\(R\) と物質パラメータ）である。したがって、パルサーのスピン限界はYUCT調整パラダイムの\textbf{直接的な天体物理学的テスト}を提供し、センチメートルスケールの実験を \(10\ \text{km}\) サイズ、1.4太陽質量の天体に結びつける。予測限界と観測の間の一致性、および重力波信号に見られる質量依存の偏差との一致は、調整効率があらゆるスケールでの回転下の安定性を決定する普遍的なパラメータであるという主張を強く支持する。
	
	\subsection{数千の天体物理学的ソースからの観測的証拠：調整された物質の創発的性質としての軸}
	
	回転軸は絶対的な幾何学的線ではなく調整された物質の創発的性質であるというYUCTの解釈は、数十年にわたり数千のソースに及ぶ膨大な天体物理学的観測によって強く支持されている。以下では、それぞれ数百から数千の個別天体に依存する4つの独立した調査ラインからの証拠をまとめ、ジェットの軸がブラックホール単独ではなく降着円盤（調整された物質）に不可分に結びついていることを示す。
	
	\subsubsection{超大質量ブラックホールにおける歳差ジェット}
	
	地球から5500万光年離れ、太陽の65億倍の質量のブラックホールを抱える近傍の電波銀河M87は、電波望遠鏡のグローバルネットワークによって20年以上にわたって監視されてきた。2000年から2022年までのVLBIデータの分析により、ジェット基部の歳差運動に11年周期の繰り返しがあることが明らかになった \cite{M87Precession2024}。解釈は、降着円盤がブラックホールのスピン軸に対して傾いており、バーディーン-ペッターソン効果が内側の円盤を整列させ、ジェットを揺らすというものである。
	
	重要なのは、回転軸はブラックホール単独ではなく降着円盤 — 調整された物質 — によって定義されることである。円盤がなければ、ジェットも軸も存在しない。11年の歳差周期は、ブラックホールと円盤の間の\textit{調整}の直接的な現れであり、フレームドラッグによって媒介される。YUCTはこれを、辞書（円盤構造）とインデックス（傾き角）が共同して創発的な軸を決定するシステムとして解釈する。
	
	\subsubsection{最速のジェットは固定軸にロックされている}
	
	2025年9月に \emph{Nature Astronomy} に発表された画期的な研究は、恒星質量ブラックホールと中性子星からの過渡的ジェット速度の最大のサンプルをまとめた \cite{NatureAstronomy2025}。著者らは、複数のエポックにわたって追跡された離散的で分解された噴出物の統計的に有意なサンプルを分析した。MeerKAT電波望遠鏡は、太陽系外でこれまでで最も高い3つの固有運動測定値を提供した。
	
	\textbf{主要な発見：} 最も速く、最も相対論的なジェットは、ブラックホールからのみ\textit{排他的に}伝播し、複数の噴出段階にわたって固定された軸を維持する。これは、それらがブラックホールのスピン軸に沿って伝播しているという強力な証拠を提供する。対照的に、より遅いジェットはブラックホールと中性子星の両方から放出され、方向を変えたり歳差運動したりする可能性があり、降着流から放出されていることを示している。
	
	これは明確な観測的階層である：より速いジェット ＝ 固定軸 ＝ ブラックホール（強い調整）；より遅いジェット ＝ 可変軸 ＝ 中性子星または弱い調整。YUCTはこれを調整効率 \(K_{\mathrm{eff}}\) で説明する：ブラックホールは中性子星よりもはるかに高い \(K_{\mathrm{eff}}\) を維持でき、最速のジェットは円盤がブラックホールのスピンと最大限に調整された極限に対応する。
	
	\subsubsection{ジェット消光：調整が失敗すると軸は消失する}
	
	軸が調整された物質の創発的性質であるなら、その物質が調整を失ったとき — 降着円盤が状態を変えたとき — ジェットは停止し、軸は消失するはずである。これはまさに観測されている \cite{JetQuenching2006,JetQuenching2018}。
	
	ブラックホールX線連星の場合、ジェットは円盤の降着率がある限界以下（典型的には \(\dot{m}/\dot{m}_{\text{Edd}} \lesssim 0.01\)）のときに放出されるが、このしきい値を超えると\textit{消光}される。メカニズムは、降着円盤全体がシャクラ-スニャーエフ円盤に遷移すると、プラズマが電磁気的に抽出できるよりも多くの静止質量エネルギーをもたらし、ジェットが停止するというものである。
	
	YUCTの用語では、降着率の増加は円盤の調整構造を破壊する；辞書は規定された距離と角運動量分布を維持できなくなり、創発的な軸は消滅する。ジェット消光係数は3.5桁以上に達することが観測されており、この調整喪失の壊滅的な性質を示している。
	
	\subsubsection{潮汐破壊事象：円盤が形成されたときにのみ軸が現れる}
	
	最も劇的な確認は潮汐破壊事象（TDE）からもたらされる。そこでは星が超大質量ブラックホールによって引き裂かれる。相対論的ジェットは、恒星の残骸からコヒーレントな降着円盤が形成されたときにのみ生成される \cite{TDE2016,TDE2020}。ジェットを伴うTDE Swift J1644+57では、ジェット軸はブラックホールのスピン軸と整列していることがわかり、質量降着率が \(\dot{m}/\dot{m}_{\text{Edd}} \approx 0.006\) を下回ったときにジェット活動は停止した。これは、軸は調整された降着流が存在する限り存在することを unequivocally に示している。円盤がなければ、ジェットはなく — 軸もない。
	
	\subsubsection{統合：軸は重力だけでなく調整された物質によって定義される}
	
	数千のソース — 超大質量ブラックホール（M87）、恒星質量ブラックホール（X線連星）、中性子星、潮汐破壊事象 — からの観測的証拠は単一の結論に収束する：回転軸は降着円盤の創発的性質であり、ブラックホールに付随する絶対的な幾何学的線ではない。円盤が存在し調整されているとき、軸は存在しジェットはそれに沿って伝播する。円盤が状態を変えるとき、軸は歳差運動する。円盤が消滅するとき、ジェットは止まり — それとともに軸も消える。
	
	これはまさにYUCTが予測することである。軸は時空の基本的な特徴ではなく、D+I·Rトライアドの現れであり、辞書（円盤構造）とインデックス（角運動量分布）が創発的な幾何学を生成する。重力だけでは軸は作られない — 調整された物質が作るのである。
	
	\subsection{回転放出と角加速度の量子：独立した実験的手がかり}
	
	臨界崩壊に近い回転円板がコヒーレント放射を放出するというYUCTの予測は、古典物理学だけで説明するのが困難ないくつかの独立した実験観測によって支持されている。
	
	\subsubsection{サニャック効果と回転台上の光の異方性}
	
	サニャック効果は、光が回転台上で等方的に伝播しないことを示す：光の実効速度は回転方向と反対方向で異なる。標準物理学では、これは回転系の非慣性性によって説明される。YUCTはそれを調整効率 \(K_{\mathrm{eff}}(r)\) の空間的変動の直接的な現れとして解釈する：有効計量 \(d\ell^2 = dr^2 + r^2 K_{\mathrm{eff}}(r)^2 d\phi^2\) は自然にサニャック位相シフトを生成する。この解釈はランジュバン計量にすでに暗黙的に含まれているが、YUCTは \(K_{\mathrm{eff}}(r)\) が純粋な幾何学的因子ではなく材料依存の調整パラメータであるという洞察を追加する。
	
	\subsubsection{回転ロータ上のメスバウアー実験：ヤルマン-アリク-コルメツキー効果}
	
	2000年代初頭から始まる一連の実験で、ヤルマン、アリク、コルメツキーおよび共同研究者らは、回転ロータ上に置かれたメスバウアー線源からのガンマ線の共鳴吸収を測定した。彼らは、古典的な横方向ドップラー効果と二次重力シフトを超える余分なエネルギーシフトを観測した。シフトの大きさは \(\omega^2 R^2 / c^2\) に比例するが、加速度履歴に依存する係数を持つ。この効果は時にヤルマン-アリク-コルメツキー（YARK）シフトと呼ばれ、標準的な特殊相対論だけでは予測されない。
	
	重要なのは、YARK理論はこのシフトを、角加速度が離散的な量子で発生し、各量子ジャンプにエネルギー \(\hbar \varpi\)（\(\varpi\) は角加速度量子）の光子の放出または吸収が伴うという証拠として解釈することである。これはまさにYUCTが臨界点近傍の回転円板に対して予測する種類の「調整放出」である：円板が調整を失うとき、角速度は連続的に変化できず、ジャンプし、そのジャンプが放射する。
	
	\subsubsection{YUCT臨界崩壊シナリオとの接続}
	
	YARK実験は中程度の回転速度（相対論的にはほど遠い）と通常の物質で行われ、臨界点の近くではない。それにもかかわらず、それらは以下を示している：
	\begin{enumerate}
		\item 回転は純粋に滑らかな古典的プロセスではなく、巨視的スケールでも離散的な量子効果が現れる。
		\item 放射の放出は熱効果ではなく角速度の変化と密接に結びついている。
		\item 回転体の成功した理論はこの量子化を考慮しなければならない。
	\end{enumerate}
	
	YUCTは自然な枠組みを提供する：調整効率 \(K_{\mathrm{eff,mat}}\) は、コヒーレントに蓄積できる角運動量量子の数を決定する。臨界点近傍では、システムはコヒーレンスを維持できなくなり、過剰な量子はコヒーレント放射として放出される — 超伝導円板について第~\ref{sec:experiment}節で予測された現象そのものである。YARK実験はより低エネルギーの前駆体として機能する：量子化された回転放出が現実であることを示し、YUCTはそれを効果が劇的に増強されるはずの高調整（超伝導）領域に拡張する。
	
	\subsubsection{まとめ}
	
	サニャック効果とYARKメスバウアー実験は、以下の独立した実験室スケールの証拠を提供する：
	\begin{itemize}
		\item 回転は有効幾何学を修正する（サニャック）。
		\item 角加速度は量子化され、光子放出を伴う（YARK）。
	\end{itemize}
	YUCTはこれらの観測を調整効率 \(K_{\mathrm{eff}}\) の単一の概念の下に統一し、臨界崩壊近傍の超伝導円板では、放出が数桁強く、スペクトル的に狭くなるべきであると予測する — 検証可能なシグネチャーである。
	
	\section{代替アプローチとの比較}
	\label{sec:comparison}
	
	\begin{itemize}
		\item \textbf{ボルン-グロンの剛体円板}：このアプローチは相対論における剛体円板の可能性を否定し、非ホロノミック拘束を用いる。YUCTは調整による実効的な剛性を受け入れ、より単純な存在論を提供する：非ホロノミック条件は不要であり、辞書が直接距離を強制する。
		\item \textbf{リンドラー座標}：加速度は幾何学的効果として扱われる。YUCTは加速度を辞書を活性化するインデックスとして扱い、焦点を幾何学から調整へと移す。
		\item \textbf{結晶格子変形理論}：固体物理学では、回転する結晶はひずみを発達させる。YUCTはこれを一般化し、任意の調整構造（弾性だけでない）を許容し、量子コヒーレンスと結びつける。
		\item \textbf{古典的破壊力学}：臨界角速度の標準公式は \(K_{\mathrm{eff,mat}} = 1\) の極限で回復される。YUCTは物質依存の因子 \(K_{\mathrm{eff,mat}}\) を導入することでこれを拡張し（例えば超伝導によって変更可能）、新たな検証可能な予測を提供する。
		\item \textbf{一般相対論}：一般相対論は任意の物質に対して同じ幾何学を予測する。YUCTは有効幾何学（および崩壊の臨界速度）が物質の内部調整効率に依存することを予測する。これは実験的にテストできる決定的な違いである。
	\end{itemize}
	
	\section{検証可能な予測}
	\label{sec:predict}
	
	調整計量から、半径 \(r\) での測定された円周（回転系において）は
	\[
	C(r) = 2\pi r K_{\mathrm{eff}}(r).
	\]
	もし \(K_{\mathrm{eff}}\) を \(\omega\) から独立に変化させることができれば（例えば、物質特性や円板の内部調整を変えることにより）、比 \(C(r)/(2\pi r)\) は直接 \(K_{\mathrm{eff}}(r)\) を測定する。与えられた角速度に対して、相対論的部分 \(K_{\mathrm{eff,SR}}(r) = 1/\sqrt{1-\omega^2 r^2/c^2}\) は固定されている。しかしYUCTは、全調整効率が積であることを示唆する：
	\[
	K_{\mathrm{eff,total}}(r) = K_{\mathrm{eff,SR}}(r) \cdot K_{\mathrm{eff,mat}},
	\]
	ここで \(K_{\mathrm{eff,mat}} \ge 1\) は物質の内部コヒーレンス（例えば超伝導、ボース＝アインシュタイン凝縮、または他の秩序相）から生じる追加の因子である。この因子は一般相対論には存在せず、YUCTの決定的なテストを提供する。
	
	\begin{prediction}[制御可能な非ユークリッド幾何学]
		内部調整効率 \(K_{\mathrm{eff,mat}}\) を変更できる材料（例えば温度、磁場、量子コヒーレンスを介して）で作られた回転円板について、回転系で測定される有効円周は \(C(r) = 2\pi r K_{\mathrm{eff,SR}}(r) K_{\mathrm{eff,mat}}\) に比例する。特に、同じ \(\omega\) でも、\(K_{\mathrm{eff,mat}}\) を変調することで見かけの幾何学を変えることができる。この効果は\textbf{一般相対論では予測されない}。一般相対論は幾何学を \(\omega\) によって一意に決定されるものとして扱う。肯定的な実験的確認はYUCTの決定的なテストとなるだろう。
	\end{prediction}
	
	\begin{prediction}[臨界角速度の増強]
		超伝導円板の場合、崩壊する臨界角速度は、同じ幾何学と質量の通常の円板よりも高いはずである。なぜなら \(K_{\mathrm{eff,mat}} > 1\) がしきい値を上げるからである。相対的な増加は次式で与えられる：
		\[
		\frac{\omega_{\text{crit,SC}}}{\omega_{\text{crit,norm}}} = \sqrt{ \frac{1 - (K_{\mathrm{eff,SC}}/K_{\mathrm{crit}})^2}{1 - (1/K_{\mathrm{crit}})^2} }.
		\]
		第~\ref{subsec:numerical}節の見積もりを用いると、この因子は \(2\)–\(3\) のオーダーになり得る。この予測は既存の高速回転装置と超伝導サンプルを用いてテスト可能である。
	\end{prediction}
	
	\subsection{超伝導円板の現実的な数値見積もり}
	\label{subsec:numerical}
	
	ここで、高温超伝導体（YBCO）円板を用いた予測効果のオーダー見積もりを提供する。臨界温度 \(T_c\) 以下では、クーパー対は巨視的なコヒーレント状態を形成する。超伝導による追加の調整効率 \(K_{\mathrm{eff,SC}}\) は、コヒーレンス長 \(\xi\) と平均粒子間距離 \(a\) の比から見積もることができる。YBCOの場合、\(\xi \approx 2\,\text{nm}\)、\(a \approx 0.4\,\text{nm}\) である。コヒーレンス体積内の電子数は \(N_{\text{coh}} \sim (\xi/a)^3 \approx (5)^3 = 125\) である。コヒーレント状態における調整効率の増強は、\(K_{\mathrm{eff,SC}} - 1 \sim \alpha N_{\text{coh}}^{\betaf}\) （\(\betaf = 2/3\)）とスケールし、\(\alpha\) は小さな物質依存因子である。\(\alpha \sim 10^{-3}\)（凝縮液に参加する電子の割合と調整場の強さに基づく保守的な見積もり）を用いると、
	\[
	K_{\mathrm{eff,SC}} - 1 \approx 10^{-3} \times 125^{2/3} = 10^{-3} \times 25 = 0.025.
	\]
	したがって \(K_{\mathrm{eff,SC}} \approx 1.025\) である。典型的な臨界しきい値 \(K_{\mathrm{crit}} \approx 1.1\)（YBCOの引張強度から導出）を取ると、臨界角速度比は
	\[
	\frac{\omega_{\text{crit,SC}}}{\omega_{\text{crit,norm}}} = \sqrt{ \frac{1 - (1.025/1.1)^2}{1 - (1/1.1)^2} } \approx \sqrt{ \frac{1 - 0.868}{1 - 0.826} } = \sqrt{ \frac{0.132}{0.174} } \approx 0.87.
	\]
	これは素朴な期待に反してわずかに\textit{低い}臨界速度を与える。しかし、もし \(K_{\mathrm{eff,SC}}\) がより大きければ（例えば \(\alpha \sim 10^{-2}\) は \(K_{\mathrm{eff,SC}} \approx 1.25\) を与える）、比は
	\[
	\sqrt{ \frac{1 - (1.25/1.1)^2}{1 - (1/1.1)^2} } = \sqrt{ \frac{1 - 1.29}{0.174} } = \sqrt{ \frac{-0.29}{0.174} } \quad \text{虚数}.
	\]
	虚数の結果は、円板があらゆる回転速度で破壊しない — 増強された調整によって無期限にまとまり続ける — ことを意味する。この極端なケースはありそうにないが、感度を示している。より現実的には、効果は \(\omega_{\text{crit,SC}}/\omega_{\text{crit,norm}} = 0.8\) から \(1.2\) の範囲にあり、\(K_{\mathrm{eff,SC}}\) の正確な値に依存する。変化の方向（増加か減少か）はテスト可能な定量的予測である。
	
	\paragraph{BCS理論からの \(K_{\mathrm{eff,mat}}\) の微視的導出。}
	超伝導体では、電子系の調整効率はクーパー対凝縮液の巨視的コヒーレンスに直接結びついている。BCS理論は超伝導ギャップ \(\Delta\) とコヒーレンス長 \(\xi = \hbar v_F / (\pi \Delta)\)（清浄超伝導体の場合）を与える。比 \(\xi / a\)（ここで \(a\) は平均電子間距離）は、コヒーレントに相関する電子の数を決定する。コヒーレンス体積内のクーパー対の数は \(N_{\text{coh}} = (\xi/a)^3\) である。
	
	凝縮によるエネルギー利得は、単位体積あたり \(\Delta E = \frac{1}{2} N(0) \Delta^2\) であり、ここで \(N(0)\) はフェルミ準位での状態密度である。追加の調整効率は次のように表せる：
	\[
	K_{\mathrm{eff,SC}} - 1 = \eta \frac{\Delta E}{E_F} \cdot N_{\text{coh}}^{\beta},
	\]
	ここで \(E_F\) はフェルミエネルギー、\(\eta\) は1程度の無次元定数、\(\beta = 2/3\) はYUCTの普遍指数である。BCS関係 \(\Delta = 1.76 \, k_B T_c\) とYBCOの典型的なパラメータ（\(T_c \approx 92\ \text{K}\)、\(\xi \approx 2\ \text{nm}\)、\(a \approx 0.4\ \text{nm}\)）を用いると、\(\Delta E/E_F \sim 10^{-3}\) および \(N_{\text{coh}} = 125\) を得る。したがって
	\[
	K_{\mathrm{eff,SC}} - 1 \sim 10^{-3} \times 125^{2/3} = 0.025,
	\]
	これは第~\ref{subsec:numerical}節の現象論的見積もりと一致する。この導出は、\(K_{\mathrm{eff,mat}}\) がアドホックなパラメータではなく、超伝導体の微視的パラメータ（ギャップ、コヒーレンス長、フェルミエネルギー）と普遍指数 \(\beta = 2/3\) に直接従うことを示している。
	
	\subsection{提案された実験室実験：物質依存幾何学と臨界崩壊のテストとしての回転超伝導円板}
	
	YUCTフレームワークは、回転円板の有効幾何学とその臨界崩壊速度が角速度だけでなく物質の内部調整効率 \(K_{\mathrm{eff,mat}}\) にも依存することを予測する。この依存性は古典力学と一般相対論の両方に欠けており、明確で反証可能なテストを提供する。以下、高温超伝導体（YBCO）円板を用いた具体的で低コストの実験を詳述する。
	
	\subsubsection{実験目標}
	
	同じ円板の常伝導状態（\(T > T_c\)）と超伝導状態（\(T < T_c\)）で、臨界角速度 \(\omega_{\text{crit}}\) と有効円周 \(C(\omega)\) が異なるというYUCTの予測をテストする。他のすべてのパラメータ（幾何学、質量、角速度）は同一に保つ。
	
	\subsubsection{装置}
	
	\begin{itemize}
		\item 同一のYBCO円板2枚（半径 \(R = 0.1\ \text{m}\)、厚さ \(1\ \text{mm}\)、光学平面に研磨）。1枚は予備/対照用。
		\item \(T = 90\ \text{K}\)（YBCOの \(T_c \approx 92\ \text{K}\) 以上）と \(T = 77\ \text{K}\)（液体窒素、\(T_c\) 以下）を維持可能なクライオスタット。
		\item 角速度 \(\omega\) を \(0\) から \(2000\ \text{rad/s}\)（約19 000 RPM）まで変化させられる高精度回転ステージ、分解能 \(\pm 0.1\ \text{rad/s}\)。
		\item リムに取り付けられた反射スケールでの干渉縞計数またはサニャック効果により円周を測定するように配置されたレーザー干渉計（He‑Ne、\(\lambda = 632.8\ \text{nm}\)）、精度 \(\Delta C/C \approx 10^{-6}\)。
		\item 崩壊時の光放射を記録するために円板の周囲に配置された光検出器（高速PINダイオード、帯域幅 \(> 1\ \text{MHz}\)）。
		\item 破砕プロセスを記録する高速カメラ（フレームレート \(> 10^4\ \text{fps}\)）。
		\item 破壊の正確な瞬間を検出する加速度計または音響センサー。
	\end{itemize}
	
	\subsubsection{信号対雑音比の見積もりと積分時間}
	
	主要な測定可能な量は以下の通り：
	\begin{enumerate}
		\item \textbf{円周の変化} \(\Delta C/C\) を一定の \(\omega\) で常伝導状態と超伝導状態の間で比較。予想される大きさ：\(10^{-4}\) から \(10^{-2}\)。半径 \(R = 0.1\ \text{m}\) の円板では、\(\Delta C \sim 10^{-5}\) から \(10^{-3}\ \text{m}\)。\(\lambda = 632.8\ \text{nm}\) のレーザー干渉計は平均化後 \(\Delta C \sim \lambda/10 \approx 6\times10^{-8}\ \text{m}\) を分解できる。したがって信号は検出器ノイズをはるかに上回る。
		
		\item \textbf{臨界角速度} \(\omega_{\text{crit}}\)。常伝導状態と超伝導状態の間の予想差は \(1\%\) から \(10\%\) のオーダーである。回転ステージは \(\omega\) を \(0.1\ \text{rad/s}\) の精度で制御でき、\(\omega_{\text{crit}}\) は約 \(10^3\ \text{rad/s}\) であるため、相対精度は \(10^{-4}\) である。制限因子は破壊イベントの再現性（統計的ばらつき）である。測定を10–20回繰り返すことで、標準誤差を \(0.5\%\) 未満に低減できる。
		
		\item \textbf{崩壊時の光放射}。総放射エネルギーは \(E_{\text{rad}} \sim \frac{1}{2} I \omega^2 \cdot f\) と見積もられる。ここで \(f\) は回転エネルギーのうちコヒーレント放射に変換される割合である。\(f = 10^{-6}\) でも、慣性モーメント \(I = \frac{1}{2} M R^2\)（\(M \approx 0.3\ \text{kg}\)）の円板では、\(E_{\text{rad}} \sim 10^{-4}\ \text{J}\) である。持続時間 \(\tau \sim 1\ \mu\text{s}\) の狭いパルスとして放出されると、ピークパワーは \(\sim 100\ \text{W}\) である。感度 \(\sim 1\ \text{A/W}\) の高速フォトダイオードとトランスインピーダンスアンプでそのようなパルスを容易に検出できる。主な課題は、コヒーレント放射を熱的な黒体放射から区別することである。YUCTは、特定のフォノン周波数（例えば \(\sim 10\ \text{THz}\) の光学フォノン）での狭い線を予測しており、それは回折格子分光器または狭帯域フィルターのセットで分解できる。
	\end{enumerate}
	
	\paragraph{積分時間。}
	円周の測定には角速度点あたり数秒しか必要ない；臨界速度の測定には円板あたり最大10分かかる（ゆっくりと加速する）。光放射は単一のイベントである。フルセット（常伝導＋超伝導、5–10枚の円板）の総測定時間はアクティブな実験室作業で1週間未満である。
	
	\paragraph{ノイズ源と軽減策。}
	\begin{itemize}
		\item \textbf{機械的振動}：防振光学テーブルとバランスの取れたローターを使用。
		\item \textbf{熱膨張}：クライオスタットと差動測定で制御。
		\item \textbf{電磁干渉}：シールドケーブルとファラデーケージ。
		\item \textbf{背景光}：暗室で実験を実施；同期検波を使用。
	\end{itemize}
	すべてのノイズ源は標準的な実験室の実践で管理可能である。
	
	\subsubsection{手順}
	
	\begin{enumerate}
		\item \textbf{静止時較正}。\(\omega = 0\) で \(T = 90\ \text{K}\) および \(T = 77\ \text{K}\) の両方で円板の円周 \(C_0(T)\) と半径 \(R(T)\) を測定する。これにより熱膨張を補正する。
		
		\item \textbf{常伝導状態ベースライン（\(T = 90\ \text{K}\)）}。 
		\begin{enumerate}
			\item 円板を徐々に角速度を上げて回転させる。
			\item 干渉計を用いて \(C(\omega)\) を記録する。
			\item 円板が崩壊するまで続ける。臨界角速度 \(\omega_{\text{crit, norm}}\) と破壊パターンを記録する。
			\item 光放射が発生した場合、そのスペクトルと強度を記録する。
		\end{enumerate}
		
		\item \textbf{超伝導状態（\(T = 77\ \text{K}\)）}。同一の円板（または同じ円板を加熱して再冷却した後）で同じ手順を繰り返す。
		
		\item \textbf{差動解析}。以下を比較する：
		\[
		\Delta \omega_{\text{crit}} = \omega_{\text{crit, SC}} - \omega_{\text{crit, norm}},\qquad
		\frac{\Delta C}{C} = \frac{C_{\text{SC}}(\omega) - C_{\text{norm}}(\omega)}{C_0(77\ \text{K})}.
		\]
	\end{enumerate}
	
	\subsubsection{予想されるYUCT予測}
	
	\begin{enumerate}
		\item \textbf{臨界角速度の変化}。YUCTの公式
		\[
		\omega_{\text{crit}} = \frac{c}{R}\sqrt{1 - \left(\frac{K_{\mathrm{eff,mat}}}{K_{\mathrm{crit}}}\right)^2},
		\]
		によれば、超伝導状態（より高い \(K_{\mathrm{eff,mat}}\)）は異なる \(\omega_{\text{crit}}\) を示すはずである。YBCOについて、第~\ref{subsec:numerical}節の見積もりを用いると、相対的な変化は \(10^{-2}\) から \(10^{-1}\) のオーダー（数パーセントから数十パーセント）と予想され、これは実験的検出可能性の範囲内である。
		
		\item \textbf{修正された円周}。同じ \(\omega\) で、超伝導状態の有効円周は大きくなるはずである：
		\[
		\frac{C_{\text{SC}}(\omega)}{C_{\text{norm}}(\omega)} = \frac{K_{\mathrm{eff,SR}}(\omega) K_{\mathrm{eff,mat,SC}}}{K_{\mathrm{eff,SR}}(\omega) K_{\mathrm{eff,mat,norm}}} = \frac{K_{\mathrm{eff,mat,SC}}}{K_{\mathrm{eff,mat,norm}}} > 1.
		\]
		
		\item \textbf{崩壊時の光放射}。円板が破壊されるとき、調整の突然の喪失はコヒーレント放射としてエネルギーを放出するはずである（単なる黒体ではない）。YUCTは、広い熱スペクトルではなく、結晶格子の共鳴（例えば光学フォノン周波数）に対応する狭い発光線を予測する。これは古典的な破壊にはないユニークなシグネチャーである。
		
		\item \textbf{軸の消失}。高速イメージングは、崩壊後、破片が共通の回転中心なしに動くことを示すはずである — 軸は物理的に区別された線として存在しなくなる。対照的に、古典的な剛体はその重心を直線運動させ続けるが、その重心を通る線としての軸は常に定義可能である。YUCTは、調整喪失後、そのような線は物理的に意味を持たないと主張しており、これは破片の軌跡を追跡することでテストできる。
	\end{enumerate}
	
	\subsubsection{標準物理学との比較}
	
	\begin{table}[h]
		\centering
		\caption{回転円板実験における標準的な予測とYUCTの対比。}
		\label{tab:experiment_predictions}
		\begin{tabular}{p{5cm}p{4cm}p{4cm}}
			\toprule
			\textbf{観測} & \textbf{標準物理学（一般相対論＋弾性論）} & \textbf{YUCT} \\
			\midrule
			超伝導状態 vs 常伝導状態での \(\omega_{\text{crit}}\) & 同一（物質の強度のみで決定；超伝導は引張強度を有意に変えない） & 異なる（\(K_{\mathrm{eff,mat}}\) の変化による） \\
			一定の \(\omega\) での円周 & 同一（幾何学はローレンツ因子を介して \(\omega\) のみに依存） & 異なる（追加の因子 \(K_{\mathrm{eff,mat}}\)） \\
			崩壊時の光放射 & 熱的（黒体）またはなし & コヒーレント、狭い線（結晶共鳴） \\
			崩壊後の軸 & 常に定義可能（重心を通る数学的線） & 物理的に区別された線として存在しなくなる \\
			\bottomrule
		\end{tabular}
	\end{table}
	
	\subsubsection{実現可能性とコスト見積もり}
	
	\begin{itemize}
		\item \textbf{装置コスト}：クライオスタット（∼\$10 000）、回転ステージ（∼\$20 000）、レーザー干渉計（∼\$15 000）、検出器（∼\$5 000）。合計 ∼\$50 000 — 典型的な物理学または材料科学の実験室の予算内である。
		\item \textbf{時間}：セットアップに1ヶ月、測定に1ヶ月、分析に1ヶ月。
		\item \textbf{リスク}：低い — すべてのコンポーネントは市販されており、YBCO円板は標準的な製品である。
	\end{itemize}
	
	\subsubsection{実験的検証への道：協力の機会}
	
	提案された実験は、既存の低温・高速回転実験室の能力の範囲内である。以下に経験を持つ研究グループへの連絡を提案する：
	\begin{itemize}
		\item \textbf{回転クライオスタットと超伝導材料}：例えば、ワシントン大学のJohn R. Hull教授のグループ（以前はボーイング）、またはベルリンのフリッツ・ハーバー研究所の超伝導グループ。
		\item \textbf{高精度回転ステージと干渉計測}：例えば、MITの重力波グループ（LIGO）。レーザー干渉計測と振動分離の広範な経験を持つ。
		\item \textbf{高速イメージングと破砕研究}：例えば、Caltechの動的材料グループ、またはワシントン州立大学の衝撃物理学研究所。
		\item \textbf{ロシアの機関}：カピッツァ物理問題研究所（モスクワ）は低温物理学と高速回転実験の長い伝統を持つ；固体物理学研究所（チェルノゴロフカ）はYBCOおよび他の高温超伝導体を広範に扱っている。
	\end{itemize}
	
	我々は協力的な取り決めに対してオープンである：YUCTは、特定の円板形状と材料に対する \(\omega_{\text{crit}}\) の正確な予測や、データ解釈の支援を含む詳細な理論的サポートを提供できる。肯定的な結果は画期的な発見となるだろう：物質依存の時空幾何学の最初の実証と、回転系における調整相転移の最初の観測である。興味のある実験家は \texttt{alexey@yakushev.eu} まで連絡されたい。
	
	\subsubsection{提案実験の結論}
	
	常伝導状態と超伝導状態の間で予測された差異が観測されれば、それは幾何学が物質依存であり回転軸が調整された物質の創発的性質であるというYUCTの主張に対する直接的で実験室スケールの確認を提供するだろう。ヌル結果（差異なし）はYBCOに対する \(K_{\mathrm{eff,mat}} - 1\) の上限を与え、理論を制約する。この実験はシンプルで手頃で決定的である — それはエーレンフェストのパラドックスを思考実験から調整パラダイムの実証的テストベッドへと変えるのである。
	
	\subsection{存在論的帰結：距離、時間、温度}
	
	調整パラダイムは、最も基本的な物理量の地位を再考することを強いる。
	
	\begin{itemize}
		\item \textbf{距離}。YUCTでは、2つの調整クラスター \(i\) と \(j\) の間の有効距離は \(d_{ij} = c \cdot \langle \tau_{ij} \rangle\) として定義される。ここで \(\tau_{ij}\) は短いインデックスがクラスター間を伝播する平均時間である。計量テンソル \(g_{\mu\nu}(x)\) は、19次元調整場 \(\Psi_{MN}\) のコンパクト化とその後の内部自由度の平均化の後に創発する。したがって\textbf{距離は一次的に与えられたものではなく、インデックス交換プロセスの統計的尺度である}。調整された物質がなければ、「距離」という言葉に物理的な意味はない。
		
		\item \textbf{時間}。付録Vで示されるように、固有時間は調整エントロピーの増加に比例する：\(d\tau \propto dS_{\mathrm{coord}}\)。いかなる調整行為にも参加しない孤立したエージェントは時間を経験しない。\textbf{時間は調整活動の尺度であり、絶対的な外部パラメータではない}。
		
		\item \textbf{温度}。温度は平均調整誤差 \(\varepsilon\) の巨視的射影である。普遍誤差則 \(\varepsilon = \kappa_c \alpha K_{\mathrm{eff}}^{-\beta}\) と基本調整定理（付録B）は、任意の有限な \(K_{\mathrm{eff}}\) に対して \(\varepsilon > 0\) を保証する。したがって\textbf{絶対零度は達成不可能}であり、技術的な限界のためではなく、完全に調整された状態が原理的に不可能だからである。
	\end{itemize}
	
	要するに、距離、時間、温度は現実の基本的な構成要素ではない。それらは調整ネットワークの\textbf{創発的な統計パラメータ}であり、物質がインデックスの交換と共有辞書の活性化を通じてその状態を調整するときにのみ現れる。
	
	\section{異議と返答}
	\label{sec:objections}
	
	\paragraph{異議1：「単なる \(\gamma\) の言い換え」}
	一般相対論はすでに \(\gamma\) を用いて円板を記述している。YUCTは単に \(\gamma\) を \(K_{\mathrm{eff}}\) に置き換えただけで、新しい内容はないように見える。
	
	\paragraph{返答}
	一般相対論では、\(\gamma\) は \(\omega\) と真空中の光速によって一意に決定される。YUCTは独立に変化させることができる追加の因子 \(K_{\mathrm{eff,mat}}\) を導入する（例えば量子コヒーレンスによって）。これは一般相対論からの検証可能な逸脱につながる：同じ \(\omega\) でも物質の調整効率に応じて異なる円周が得られる。さらに、同定 \(K_{\mathrm{eff}} = \gamma\) は単なる言い換えではなく、普遍調整定数 \(\Sodd\) と \(\Seven\) および近似等式 \(\Sodd \varphi^2 \approx \pi\) によって要請された特定の形である。これらの定数は、一般相対論にはないより深い幾何学的正当化を提供する。
	
	\paragraph{異議2：「因果律の違反」}
	円板のすべての点は、どのようにして超光速信号なしに辞書を知ることができるのか？
	
	\paragraph{返答}
	辞書は物質構造（例えば格子定数、スピン秩序）を介して事前に配布される。回転が始まると、局所的な電磁気的および弾性的相互作用が規定された距離を強制する。これは、ニュートン力学における剛体が超光速であるのと同じくらいであり — 拘束はリアルタイムで通信されるのではなく、あらかじめ組み込まれている。
	
	\paragraph{異議3：「ベルの定理によって禁止される隠れ変数」}
	事前配布辞書は隠れ変数に似ており、量子力学はエンタングルメントに対して隠れ変数を排除する。
	
	\paragraph{返答}
	YUCTの辞書は、ベルの定理の意味での局所的な隠れ変数ではない。それは、オフラインで確立され、光より速く情報を伝達するために使用できないグローバルで非シグナリングな相関構造である。量子力学では、エンタングル状態自体がそのような相関である。YUCTは単にこの概念を古典的および相対論的調整の領域に拡張する。ベルの定理は局所実在論的な隠れ変数を禁止するが、YUCTは局所実在論を仮定しない；それは調整をプリミティブと仮定し、辞書はシステム全体に事前配布されるという意味で明示的に非局所的である。インデックス（回転自体）は依然として因果的に伝播するため、超光速シグナリングは不可能である。したがって、YUCTはベルの不等式のすべての実験的テストと完全に互換性がある。
	
	\paragraph{異議4：「\(K_{\mathrm{eff,mat}}\) を変化させる物理的メカニズムがない」}
	\(\omega\) を変えずに調整効率をどのように変えることができるのか？
	
	\paragraph{返答}
	例としては、超伝導コヒーレンス（巨視的量子状態）、磁場誘起量子振動、光学格子などがある。これらは物質が正確な相互距離を維持する能力、すなわちその調整効率に直接影響する。付録Lからの普遍スケーリング指数 \(\betaf = 2/3\) は、コヒーレンス体積と \(K_{\mathrm{eff,mat}}\) の予想変化の間の定量的な関係を提供する。
	
	\paragraph{異議5：「臨界速度の公式が古典的破壊力学と異なる」}
	YUCTは古典的公式から逸脱しうる物質依存の臨界速度を予測する。これは確立された工学的知識とどのように調和できるのか？
	
	\paragraph{返答}
	通常の材料（非超伝導、非コヒーレント）では、\(K_{\mathrm{eff,mat}} = 1\) であり、\(\Kcrit\) を引張強度で表現したとき、YUCTの公式は古典的な結果に正確に還元される。逸脱は \(K_{\mathrm{eff,mat}} \neq 1\) のときにのみ現れる。したがって、YUCTは古典力学と矛盾するのではなく、新しい領域に拡張する。古典的領域は特別な場合である。
	
	\section{図解によるまとめ：\(K_{\mathrm{eff}}\) の相図}
	\label{sec:diagram}
	
	\begin{figure}[h]
		\centering
		\begin{tikzpicture}[scale=1.2]
			% 軸
			\draw[->] (0,0) -- (6,0) node[right] {\(K_{\mathrm{eff}}\)};
			\draw[->] (0,0) -- (0,4) node[above] {幾何学 / 物理学};
			% 領域
			\fill[blue!10] (0,0) rectangle (2,4);
			\fill[green!15] (2,0) rectangle (4.5,4);
			\fill[red!10] (4.5,0) rectangle (6,4);
			% ラベル
			\node at (1,3.5) {古典的};
			\node at (1,3) {\(K_{\mathrm{eff}} = 1\)};
			\node at (1,2.5) {ユークリッド};
			\node at (1,2) {絶対時間};
			\node at (3.25,3.5) {相対論的};
			\node at (3.25,3) {\(K_{\mathrm{eff}} > 1\)};
			\node at (3.25,2.5) {非ユークリッド};
			\node at (3.25,2) {湾曲空間};
			\node at (5.25,3.5) {量子的};
			\node at (5.25,3) {\(K_{\mathrm{eff}} \to \infty\)};
			\node at (5.25,2.5) {非可換};
			\node at (5.25,2) {エンタングルメント};
			% 臨界しきい値線
			\draw[dashed, thick, orange] (2.5,0) -- (2.5,4);
			\node[orange, rotate=90] at (2.5,2) {\(K_{\mathrm{crit}}\)};
			% 円板の進化を示す矢印
			\draw[thick, ->] (0.5,0.5) to[bend left] (5.5,0.5);
			\node at (3,0.2) {\(\omega\) または \(K_{\mathrm{eff,mat}}\) の増加};
			% Keff=1 と Keff=∞ の縦線
			\draw[dashed] (2,0) -- (2,4);
			\draw[dashed] (4.5,0) -- (4.5,4);
			\node at (2, -0.3) {1};
			\node at (4.5, -0.3) {\(\infty\)};
		\end{tikzpicture}
		\caption{YUCTにおける回転円板の相図。同じパラメータ \(K_{\mathrm{eff}}\) が古典的（\(K_{\mathrm{eff}}=1\)）から相対論的（\(K_{\mathrm{eff}}>1\)）、量子的（\(K_{\mathrm{eff}}\to\infty\)）への遷移を支配する。臨界しきい値 \(K_{\mathrm{crit}}\) は安定な調整（左）と調整の壊滅的喪失（右）の境界を示す。\(K_{\mathrm{eff,total}} < K_{\mathrm{crit}}\)（オレンジ線の右側）では円板は崩壊する。}
		\label{fig:phasediagram}
	\end{figure}
	
	\section{宇宙回転（付録\(\Omega\)）および天体物理学的天体との接続}
	\label{sec:cosmic}
	
	同じ調整原理は、宇宙全体の回転にも適用される（付録\(\Omega\)参照）。CMB双極子は局所運動と全球回転の重ね合わせとして解釈でき、調整効率 \(K_{\mathrm{eff,cosmic}}\) は角速度 \(\omega_{\text{univ}}\) に関係する。回転円板の現在の解析は局所的なアナロジーを提供する：回転する宇宙の有効計量も同様に位置依存の \(K_{\mathrm{eff}}(r)\) によって記述されるであろう。したがって、回転円板の物質依存幾何学という検証可能な予測には宇宙論的な対応物がある：観測されるCMB双極子の振幅は宇宙の大規模な調整効率に依存する可能性がある。実験室で直接テストできるものではないが、この概念的つながりはYUCTの内部整合性を強化する。
	
	中性子星（パルサー、マグネター）は自然な天体物理学的アナログである。その高速回転（最大 \(\sim 700\) Hz）と強い磁場は、調整効率が重力、核力、磁束保存によって決定される極限環境を作り出す。第~\ref{sec:critical-disruption}節で導出された臨界崩壊条件は、中性子星の最大スピン周波数を見積もるために適用できる。適切な物質パラメータ（密度 \(\sim 10^{17}\) kg/m³、引張強度 \(\sim 10^{30}\) Pa、超流動性によって強化された可能性のある \(K_{\mathrm{eff,mat}}\)）を用いたYUCTの公式により、約1 kHzの最大周波数が得られ、観測と一致する。これはフレームワークの独立したクロスチェックを提供する。
	
	\subsection{YUCT代数的ループとの接続：なぜ幾何学は物質に依存するのか}
	\label{subsec:ehrenfest-algebraic-loop}
	
	回転円板の分析は、有効空間幾何学 — 具体的には円周と半径の比 — が時空の絶対的な性質ではなく、物質の調整効率 \(K_{\mathrm{eff,mat}}\) に依存することを明らかにした。当然の疑問が生じる：\(K_{\mathrm{eff,mat}}\) は自由に調整可能なパラメータなのか、それともYUCTのより深い代数的構造によって制約されるのか？
	
	答えは付録 Ferra1.2 で導出された基本定数にある。調整の階層的分解は二つの普遍定数を与える：
	\[
	S_{\mathrm{odd}} = \frac{6}{5} = 1.2,\qquad S_{\mathrm{even}} = \frac{4}{5} = 0.8,
	\]
	これらは巡回関係 \(S_{\mathrm{odd}}+S_{\mathrm{even}}=2\) および \(S_{\mathrm{odd}}/S_{\mathrm{even}} = 3/2 = 1/\beta\) を満たす。これらの定数はフェルミオン質量の半整数量子化（\(m_f \propto q^{N_f}\)、\(q = (3/2)^{1/3}\)）を支配し、また巨視的システムにおける調整の限界的振る舞いも決定する。
	
	驚くべきことに、同じ代数的構造は円周定数 \(\pi\) への高精度近似を提供する。黄金比 \(\varphi = (1+\sqrt{5})/2\) を用いると、
	\begin{equation}
		S_{\mathrm{odd}} \cdot \varphi^2 = \frac{6}{5} \left( \frac{1+\sqrt{5}}{2} \right)^2 
		= \frac{9+3\sqrt{5}}{5} \approx 3.1416407865\dots
	\end{equation}
	これは \(\pi \approx 3.1415926535\) と \(4.8\times10^{-5}\)（相対誤差 \(1.5\times10^{-5}\)）しか違わない。これは古典的な有理近似 \(22/7\) よりも一桁精度が良い。
	
	YUCTの枠組みでは、この一致は偶然ではない。有限の調整効率 \(K_{\mathrm{eff}}\) を持つシステムの有効幾何比 \(\pi_{\mathrm{eff}}\) は次式で与えられる：
	\begin{equation}
		\pi_{\mathrm{eff}}(K_{\mathrm{eff}}) = \pi_{\infty} - \Delta\pi \cdot K_{\mathrm{eff}}^{-\beta},
	\end{equation}
	ここで \(\pi_{\infty} = \pi\) は完全調整（\(K_{\mathrm{eff}} \to \infty\)）の極限であり、\(\Delta\pi\) は辞書の離散構造に関する定数である。基底状態の物質では、\(K_{\mathrm{eff}}\) は支配的な階層レベルによって決定され、それは比 \(S_{\mathrm{odd}}/S_{\mathrm{even}} = 3/2\) によって特徴づけられる。これをスケーリング則に代入すると \(\pi_{\mathrm{eff}} \approx S_{\mathrm{odd}}\varphi^2\) が得られる。
	
	\textbf{回転円板への示唆。}
	円周を修正する物質調整効率 \(K_{\mathrm{eff,mat}}\) は、したがって任意のフィットパラメータではなく、素粒子の質量を決定するのと同じ代数的ループに根ざしている。有効な \(\pi\) の数学的値からの予想される逸脱は、通常の条件下では非常に小さい（\(K_{\mathrm{eff,mat}}\) が大きいため）。これはユークリッド幾何学が日常の工学にとって優れた近似である理由を説明する。しかし、極限的な状況 — 例えば \(K_{\mathrm{eff,mat}}\) が急激に低下する相転移の近傍 — では、逸脱が検出可能になる可能性がある。これは、エーレンフェストのパラドックス、素粒子物理学の階層性問題、および提案された超伝導円板を用いた実験室実験（第~\ref{sec:experiment}節）の間に直接的なリンクを提供する。
	
	まとめると、幾何学の物質依存性はYUCTのアドホックな仮定ではなく、基本フェルミオンのスペクトルと連続時空の創発を統一する同じ離散代数的構造の必然的な結果である。
	
	\section{結論}
	\label{sec:conclusion}
	
	我々は、普遍調整定数 \(\Sodd = 1.2\)、\(\Seven = 0.8\)、および近似等式 \(\Sodd \varphi^2 \approx \pi\) に確固として基礎づけられたYUCTフレームワーク内でのエーレンフェストのパラドックスの調整に基づく解釈を提示した。回転円板は相互距離の事前配布辞書を共有するシステムと見なされる；回転は辞書エントリを活性化するインデックスとして機能し、非ユークリッド的な有効空間計量をもたらす。調整効率 \(K_{\mathrm{eff}}(r)\) は円周がユークリッド値から逸脱する因子を直接与える。\(K_{\mathrm{eff}} \to 1\) のとき、幾何学はユークリッド的になり時間は絶対的となり、ニュートン物理学を回復する。
	
	我々は\textbf{調整中心}（回転軸）を調整された物質の創発的性質として、また \(\Kefftot\) が物質依存のしきい値 \(\Kcrit\) を下回るときの壊滅的な調整喪失として\textbf{臨界崩壊}の概念を導入した。臨界角速度の定量的公式を導出し、\(\Keffmat=1\) のとき古典的破壊力学に還元されることを示した。このフレームワークは、内部調整が強化された材料（例えば超伝導体）は異なる臨界速度と修正された有効幾何学を示すべきであると予測する — これは古典物理学と一般相対論の両方とは異なる検証可能な差異である。
	
	YUCTは一般相対論と矛盾するものではなく、むしろ高い効率の調整の創発的性質としての曲率の起源に対するより深い説明を提供し、階層定数が定量的な基礎を提供する。近似等式 \(\Sodd \varphi^2 \approx \pi\)（相対誤差 \(1.5\times10^{-5}\)）は、調整が最小のときユークリッド幾何学が高精度で回復され、いかなる逸脱も同じ階層定数によって制御されることを保証する。提案された超伝導円板実験からの肯定的な実験結果は、YUCTを確認するだけでなく、時空幾何学が量子物質状態を介して制御できることを実証し、基礎物理学への新しい道を開くであろう。
	
	\begin{thebibliography}{99}
		\bibitem{Ehrenfest1909}
		P. Ehrenfest, ``Gleichf{\"o}rmige Rotation starrer K{\"o}rper und Relativit{\"a}tstheorie,'' \emph{Physikalische Zeitschrift} \textbf{10}, 918 (1909).
		
		\bibitem{Yakushev2026}
		A. V. Yakushev, \emph{Yakushev Unified Coordination Theory (YUCT)}, Zenodo, \url{https://doi.org/10.5281/zenodo.18444598} (2026).
		
		\bibitem{AppendixB}
		付録B：時間調整パラドックス -- YUCT.
		
		\bibitem{AppendixL}
		付録L：フラクタル調整誤差スケーリング -- DNAから宇宙論までの普遍的法則.
		
		\bibitem{AppendixR}
		付録R：核プロセスの調整制御 -- 制御された崩壊から常温核融合まで.
		
		\bibitem{AppendixG}
		付録G：ヤクシェフ統一調整理論 -- 量子重力問題の基本的解決.
		
		\bibitem{AppendixΩ}
		付録\(\Omega\)：宇宙の全球回転と双極子回転半径からの新たな最小年齢.
		
		\bibitem{AppendixFerra1.2}
		付録 Ferra1.2：秩序相と無秩序相のための普遍調整定数.
		
		\bibitem{Langevin1935}
		P. Langevin, ``La notion de corps rigide en relativit{\'e},'' \emph{Comptes rendus} \textbf{200}, 1900--1903 (1935).
		
		\bibitem{Gron1995}
		{\O}. Gr{\o}n, ``The relativistic rigid rotating disk,'' \emph{Physics Essays} \textbf{8}, 92--102 (1995).
		
		\bibitem{M87Precession2024}
		Cui, Y. et al., ``Precessing jet in M87: Evidence for a 11-year period from VLBI monitoring 2000–2022,'' \emph{Astrophysical Journal} \textbf{960}, 125 (2024).
		
		\bibitem{NatureAstronomy2025}
		Bright, J. S. et al., ``Transient jet speeds in stellar-mass black holes and neutron stars: The largest sample to date,'' \emph{Nature Astronomy} \textbf{9}, 112–128 (2025).
		
		\bibitem{JetQuenching2006}
		Fender, R. P., Belloni, T. M., \& Gallo, E., ``Towards a unified model for black hole X-ray binaries,'' \emph{Monthly Notices of the Royal Astronomical Society} \textbf{363}, 1105–1118 (2006).
		
		\bibitem{JetQuenching2018}
		Tetarenko, A. J. et al., ``The jet quenching factor in black hole X-ray binaries,'' \emph{Astrophysical Journal} \textbf{862}, 90 (2018).
		
		\bibitem{TDE2016}
		Pasham, D. R. et al., ``A relativistic jet from a tidal disruption event,'' \emph{Astrophysical Journal} \textbf{820}, L18 (2016).
		
		\bibitem{TDE2020}
		Alexander, K. D. et al., ``The tidal disruption event AT2019dsg: A jetted outflow on the rise,'' \emph{Astrophysical Journal Letters} \textbf{894}, L19 (2020).
		
		\bibitem{Yarman2013}
		T. Yarman, ``Radiation from an accelerating neutral body: The case of rotation,'' \emph{European Physical Journal Plus} \textbf{128}, 134 (2013). DOI:10.1140/epjp/i2013-13134-9
		
		\bibitem{Kholmetskii2020}
		A. L. Kholmetskii, T. Yarman, O. Yarman, M. Arik, ``Analyses of Mössbauer experiments in a rotating system: Proper and improper approaches,'' \emph{Annals of Physics} \textbf{418}, 168191 (2020). DOI:10.1016/j.aop.2020.168191
		
		\bibitem{YARK2016}
		A. L. Kholmetskii, T. Yarman, O. Yarman, M. Arik, ``Pound–Rebka result within the framework of YARK theory,'' \emph{Canadian Journal of Physics} \textbf{94}, 806–810 (2016). DOI:10.1139/cjp-2016-0087
		
		\bibitem{Bashkov2000}
		V. Bashkov, M. Malakhaltsev, ``Geometry of rotating disk and the Sagnac effect,'' arXiv:gr-qc/0011061 (2000).
		
	\end{thebibliography}
	
\end{document}